% Français 12357, batch 3 — Gallica views 41–56 (the last batch; sixteen views).
% Pass: Opus 5.5 (claude-opus-5-5), 2026-09-22 — first pass, unchecked against the views by a human.
% Read from the local mirror only (archives/fr-12357/f0041.jpg … f0056.jpg):
% 1000-px thumbnails, then four full-width bands per written view at 2400 px,
% then tight crops of 1800 px for doubtful words, the tables and the magic
% squares. Each view was drafted to archives/scratch/fr-12357-b3/ as it was
% read and the file assembled from the drafts. No edition of Mersenne, no
% edition of the Harmonie universelle and no other transcription was opened:
% every reference to the book below is copied as the leaf gives it and none
% has been checked against the book.
%
% What the views are. Greyscale scans, portrait, 3919 × 5747 (3920 × 5748 for
% v. 56), one page per view: a left page (verso) and a right page (recto)
% alternate, the edge of the facing page and its show-through visible at the
% gutter. On versos the ends of lines run into the gutter and are often
% hidden; they are completed in \add{} only when the word is certain, and
% otherwise marked \ill{} with a note. Several pages (v. 43, 44, 45, 50, 51)
% carry notes, tables or figures in the inner margin that appear again, from
% the other side of the gutter, on the facing view; each is transcribed once,
% on the page it is written on, and the notes say so.
%
% The hand. One copyist throughout the sixteen views (the abbreviations
% match those batch 2 records; the earlier views were not compared): a small, fast, upright seventeenth-century French
% cursive in dark ink, with a more careful and larger module for the Latin
% piece of v. 53–55. Long s is set s. The copyist's abbreviations are kept
% as written and expanded, where needed, only in notes: pō (pour), co^e
% (comme), f^e (faire), p^r / p^re (premier, première), so^e (somme), p^ties
% (parties), lesd^t. / desd^t. / aud^t. (lesdits…), commencem^t, tous^rs
% (tousjours), m^r, nō (nous), -aõn (-ation: observaõn, multiplicaõn,
% Denominaõn), cæt. / caet. (cætera); in the French the word « est » is a
% siglum written as a small raised loop and is set « est » silently. In
% Latin: the -ur ending as a z-shaped sign (expanded silently, as in the
% Roberval volume), -bus as a superscript 9 (kept: oĩb^9), -que as q; or q
% with a hook (expanded), the suspended nasal as a tilde (kept where the
% expansion is not certain: coĩ, nõ, oẽs, viã), id est as « id ẽ ».
% « i » for « j » and « u »/« v » stand as written.
%
% Notation. Superscripts written by the copyist are \textsuperscript{}. The
% sign like a figure 4 that stands for the natural (b-quarre) in the tables
% of v. 43 and v. 55 is set \natural, as in batch 2. A note letter followed
% by a small x-shaped stroke in the table of v. 43 (Cx, Dx, Fx, Gx) is set
% with an x; that it marks the sharp is the pass's inference. The copyist
% draws small arcs over pairs of table columns and under the intervals of a
% series (v. 43, 51); these are described, not drawn. Magic squares, tables
% of aliquot parts and dice tables are set as arrays, cell by cell, exactly as
% read; the checks made on them (row, column and diagonal sums) are the pass's
% own and are stated in the notes.
%
% Blank or text-less views, absent from the file: view 56, the verso of the
% last written leaf, blank but for show-through of v. 55 and the round stamp
% « BIBLIOTHEQUE IMPERIALE MSS. ». View 52 has no text of its own (it is
% v. 51 seen through the paper) but keeps a \page{} for a note, because one
% pale line at its foot reads the right way round and repeats the last line
% of v. 50; it is described, not transcribed.
%
% What the sixteen views hold. Two kinds of copy: remarks tied to a page of
% the book, headed « p. N. lig. N. », « au commencem^t », « coroll. N »,
% « observation N » (v. 41–44), then, from the rubric « Au papier blanc
% apres tous les livres de l'harmonie » (v. 45), the notes written on the
% blank leaves at the back of the book (v. 45–55), which cite the book only
% occasionally.
%
%   v. 41      Continues v. 40 mid-sentence (« et que d'un / tuyau de 12
%              pieds »): the speed of water out of a pipe against that of an
%              arquebus ball (43200 pieds; 100 toises in a second); water
%              passing from a wide pipe into a narrow one. Then remarks on
%              p. 6 coroll. 1 line 2 and line 13 (a lead weight 1890 times
%              heavier than a bladder), p. 7 line 5, p. 7 coroll. 5 (against
%              Galileo: lead should fall 42 feet in water in 2'' and falls
%              only 12), p. 9 line 25, « Ibid. coroll. 8 » (dated in the text
%              « le 4^e Juin 1638 »), and p. 10 « seconde observation »: the
%              motion compounded of natural and violent that describes a
%              parabola ABC, with a figure in the margin, and the correction
%              « Il faut corriger tout cecy » (arrows show the descent slower
%              than the ascent).
%   v. 42      p. 13 line 25 (testing it with a crossbow bolt from a tower);
%              p. 15 « observaõn 4^me » line 7 (whistling into a glass, a
%              cousin of Huygens de Zuylichem, news from Utrecht); p. 16
%              before the 5th observation (Galileo's drops that split when the
%              glass leaps to the octave); the 5th observation (Descartes's
%              letters); p. 17 « au commencem^t »: a refutation in five points
%              of Beaugrand's geostatics (the balance whose arms point to the
%              centre of the earth), with a figure in the margin; p. 18, the
%              6th observation.
%   v. 43      p. 22, 9th observation: a large table of the thirds, fourths,
%              fifths and sixths of the monochord in four columns (perfect,
%              « aprochants », imperfect, « du tout esloignez ») against the
%              equal-tempered « Accord de Ga[llé] » by mean proportionals, on
%              a string of 100000, with the notes that agree and disagree;
%              then a paragraph on organ-builders' ears, 53 commas to the
%              octave, and the equal semitone. It runs on to v. 44.
%   v. 44      End of that paragraph; p. 25 « au commencem^t » (the roulette:
%              its area for any ratio, its tangents, the solid of revolution
%              8 to 5, centres of gravity, « voyez les lettres de m^r
%              descartes »); p. 26 « au commencem^t », in Latin: Descartes's
%              rule for amicable numbers (284/220, 18416/17296,
%              9437056/9363584); the 13th observation line 4; multiply perfect
%              numbers (lists of double, triple, quadruple; « trouvé par m^r
%              descartes »); a boxed note on the quintuple and septuple.
%   v. 45      « Au papier blanc apres tous les livres de l'harmonie »: a magic
%              square of 22 containing those of 16, 12 and 10 (rows 5335),
%              with marginal corrections of the sums given « selon le pere
%              mersenne »; the planetary squares of 9 (Moon) and 8 (Mercury),
%              a square of 4; the variety of the magic squares (880 ways for
%              4; 6742112993280000 for the square of 9).
%   v. 46      The square of 14 in « une lettre de m^r. fermat au pere mersenne
%              du p^r [?] Apvril 1640 »; another square of 4; varieties of the
%              squares of 14 and 22 (three very large numbers); the challenge
%              « Arithmeticis proponitur numerus 2027651281 »; rules 1 and 2
%              for the number of aliquot parts of a number and for the least
%              number with a given number of aliquot parts (108 → 11; 11 → 72).
%   v. 47–48   Continuation: numbers of aliquot parts (20, 360, 400, 5000,
%              15000), rule 3 (the number whose aliquot parts sum to a given
%              number), rule 4 (the sum of the aliquot parts: 20, 5040,
%              560400, 666, 12000, 5000, 888), triples of numbers in
%              continued proportion contained in a number (91, 49, 24843),
%              tests of non-primality (divisibility by 5, 3, 9), and 6000
%              « nombre des années de la durée du monde » (39 aliquot parts,
%              their sum 13344), announced « a la page suyvante ».
%   v. 49      A loose slip lying over v. 50: the throws of two dice, sums
%              2 to 12, first written with all pairs up to 10+1 and then the
%              impossible ones struck out.
%   v. 50      The 39 aliquot parts of 6000 (table); 6679 = 23 · 29; 1373
%              prime; whether 1073602561 is prime (divide by primes up to
%              8191); in Latin: sums of triangular numbers, the civic number
%              of Plato 5040 (59 aliquot parts, their table and the
%              decomposition 2·2·2·2·3·3·5·7), five problems of stones set out
%              along a road (arithmetical and geometrical spacing), sums of
%              cubes and of squares, running on to v. 51.
%   v. 51      Sum of the first 30 squares (9455); « Quarrez quarrez ou 4^mes
%              puissances et leurs differences » by centred hexagonal numbers;
%              fifth powers; « Des differens ieux des dez »: three dice (216,
%              the counts 1, 3, 6, 10, 15, 21, 25, 27), two dice; the first
%              fourth powers and their differences. Two dice tables in the
%              margin.
%   v. 53–55   A separate piece in Latin: « Academia Parisiensis viros
%              Clarissimos Galilaei familiares et amicos lyncaeos precatur, uti
%              sequentibus in dialogorum libros notis respondeant »: fifteen
%              numbered objections to Galileo's dialogues, by page (2–3, 23,
%              17, 28, 42, 50, 54, 70, 73, 103, 166, 236, 288, 280), with a
%              figure of a gun and a falling ball (v. 54); a request for the
%              weight and size of the flasks with which Galileo weighed air and
%              for his treatise on percussion; then « Haec sunt quae ipsi
%              Galilaeo etiamnum viventi … scripserant quorum missionem tanti
%              viri obitus praevenit », asking his friends to send their answers
%              to « D. Santininum » in Genoa, dated « Lutetiae Calendis Julii
%              anni 1643 ». Below it, on v. 55, a diagram of the seven Greek
%              modes as staircased scales (Hypodorien to Mixolidien) with a
%              French legend citing « m^r Doni ».
%
% References to the book, exactly as written on the leaves. v. 41: « p. 6.
% coroll. 1^er ligne 2. », « lig. 13. », « p. 7 lig. 5. », « p. 7. coroll.
% 5 », « p. 9. lig. 25. », « Ibid. coroll. 8. », « la 29 et 32. prop. de
% l'hydraulique pneumatique », « p. 10. seconde observation », « voyez la
% balistique ». v. 42: « p. 13. lig. 25. », « p. 15. observaõn 4^me. lig.
% 7. », « p. 16. avant la 5^me observation », « 5^e observation », « p. 17.
% au commencem^t. », « la 6^me. raison p. 61 de l'utilité de l'harmonie prop.
% 18. », « p. 18. en la 6^e observation », « la 6^me proposition du 5^me livre
% des instruments ». v. 43: « p. 22. observation 9^me ». v. 44: « p. 25. au
% commencem^t », « p. 26. au commencem^t. », « xiij observation. lig. 4. »,
% « la preface des hydrauliques ». v. 45: « Au papier blanc apres tous les
% livres de l'harmonie ». v. 46: « la table de la 108 page du livre de la
% voix ». The leaves never name the part of the book that pages 6 to 26
% belong to; that part is arranged in numbered « observations » and
% « corollaires », which is all the pass can say. References to other works
% are also copied as written: Galileo's « nouvelles pensées » page 67 after
% the 15th article (v. 41) and « la preface sur les nouvelles pensées de
% Galilée » (v. 44), « les reflexions physicomathematiques \ill{} chap. 22 »
% (v. 41), « le livre de m^r de la brosse » and Descartes's letters (v. 42,
% 44), a letter of Fermat to Mersenne of April 1640 (v. 46), and the pages of
% Galileo's dialogues in the Latin piece (v. 53–54).
%
% Numbers on the leaves. (1) An ink series, one figure per leaf at the top
% right of each recto only, in a finer stroke than the text: 20 (v. 42), 21
% (v. 44), 22 (v. 46), 23 (v. 48), 24 (v. 51), 25 (v. 53), 26 (v. 55, the
% figure underlined). No verso carries a number. It ends at 26 on the last
% written leaf, which is the catalogue's count of 26 feuillets, so it behaves
% like a foliation; but it is in ink, not the thin grey pencil this project
% takes for the library's foliation, and batch 2 of this volume records the
% same series in notes without \folio{}. This batch follows it: every figure
% is recorded in the view's note and no \folio{} is written. A human look at
% the leaves would settle whether it is the library's count. (2) A small
% isolated « 3 » at the foot of v. 42, bottom right, in another pen — perhaps
% a gathering mark (inference). (3) A small « 1 » or « i » at the top left of
% v. 54. (4) Every other number on these pages is content (page references to
% the book, the numbers of the calculations and tables), including the « 154 »
% written in the margin of v. 46 opposite a blotted cell of the square of 22
% on v. 45, and the numbers 1–12 of the dice table of v. 51 seen on the edge
% of v. 50. None of the numbers is the copyist's own pagination.
%
% Across the edges of the batch. View 41 opens mid-sentence: view 40 ends
% « si l'eau monte d'autant plus viste perpendiculairement qu'elle sort plus
% viste, et que d'un », and view 41 continues « tuyau de 12 pieds de
% hauteur… ». Within the batch, each page ends on a catchword or a broken
% sentence taken up on the next written page: 41→42 (« p. 13 lig. 25. »),
% 43→44 (« que le Iugem^t »), 45→46 (« voyez le quarré »), 46→47, 47→48,
% 48→50 (« a la page suyvante »), 50→51 (« 10, supersunt »), 54→55 (« aeris
% pondus »). The batch ends at the end of the volume.
% The planetary signs of v. 45, which the PDF font lacks, are set as their
% French names in brackets (\add{Lune}, \add{Mercure}…), the sun as $\odot$.

\documentclass[11pt,a4paper]{article}
% The preamble shared by every transcription of the mathematicians' manuscripts in Gallica.
%
% It rests on one idea: **uncertainty compounds, it does not resolve.** A
% transcription that smooths over a doubtful word has destroyed the only thing
% it was made for — a reader must be able to tell, at every point, what was on
% the page and what was supplied. The apparatus macros below are therefore the
% critical apparatus, not decoration.
%
% The same commands are recognised by `scripts/render.mjs`, which produces the
% site's reading view, and by `scripts/tei.mjs`. A macro added here must be
% added there too, or rendering fails loudly — which is the wanted behaviour.
%
% Comments here are English, like the rest of the repository. The documents
% this preamble sets are French, and that is deliberate: see the skills.

\ProvidesPackage{germain}[2026/09/15 Transcriptions of mathematicians' manuscripts in Gallica]

\usepackage{amsmath,amssymb,amsthm}
\usepackage{mathrsfs}
% Kept from the parent preamble so the renderer's diagram support stays
% available; these manuscripts draw no commutative diagrams, and nothing here uses it.
\usepackage{tikz-cd}
\usepackage[english,french]{babel}
\usepackage{fontspec}

% --- Greek ------------------------------------------------------------------
% The text face stays fontspec's default, Latin Modern, which has no Greek:
% XeTeX drops every glyph it lacks with a line in the log and nothing on the
% page, and the Greek words the manuscripts quote vanished from the PDFs. So
% the Greek and Greek Extended blocks get a XeTeX character class of their own,
% and entering or leaving a run of it switches the family to CMU Serif — the
% same Computer Modern design, with polytonic Greek — and back. Only the
% family changes: a Greek word in \textit or \textbf keeps its shape and series.
% The transitions to and from every other class carry the tokens that class
% has towards an ordinary letter, so French spacing before « ; » or inside
% « » still holds after a Greek word. No grouping, so no brace or \endgroup
% can come out unbalanced; maths is left alone, since XeTeX inserts nothing
% there. Kept identical in `exercices.sty`.
\makeatletter
\newfontfamily\germainGreekfont{cmun}[
  Extension=.otf, NFSSFamily=germaingreek,
  UprightFont=*rm, ItalicFont=*ti, SlantedFont=*sl,
  BoldFont=*bx, BoldItalicFont=*bi, BoldSlantedFont=*bl]
\newXeTeXintercharclass\germain@greek
\def\germain@greekrange#1#2{%
  \@tempcnta=#1\relax
  \loop
    \XeTeXcharclass\@tempcnta=\germain@greek
    \ifnum\@tempcnta<#2\advance\@tempcnta\@ne\repeat}
\germain@greekrange{"0370}{"03FF}
\germain@greekrange{"1F00}{"1FFF}
\def\germain@greekprev{\rmdefault}
\def\germain@greekin{%
  \edef\germain@greekprev{\f@family}\fontfamily{germaingreek}\selectfont}
\def\germain@greekout{\fontfamily{\germain@greekprev}\selectfont}
\def\germain@greekjoin{%
  \XeTeXinterchartokenstate=\@ne
  \@tempcnta=\z@
  \loop
    \ifnum\@tempcnta=\germain@greek\else
      \XeTeXinterchartoks\germain@greek\@tempcnta=\expandafter{%
        \expandafter\germain@greekout\the\XeTeXinterchartoks\z@\@tempcnta}%
      \XeTeXinterchartoks\@tempcnta\germain@greek=\expandafter{%
        \the\XeTeXinterchartoks\@tempcnta\z@\germain@greekin}%
    \fi
    \ifnum\@tempcnta<4095\advance\@tempcnta\@ne\repeat}
% babel-french sets its own classes, and saves and restores the interchar
% state, each time a language is selected: join again after it does.
\addto\extrasfrench{\germain@greekjoin}
\addto\extrasenglish{\germain@greekjoin}
\AtBeginDocument{\germain@greekjoin}
\makeatother

\usepackage{geometry}
\geometry{a4paper,margin=25mm,marginparwidth=28mm}
\usepackage{marginnote}
\usepackage{xcolor}
\usepackage[normalem]{ulem}
\setlength{\emergencystretch}{3em}
\usepackage{hyperref}
\hypersetup{colorlinks=true,linkcolor=black,urlcolor=black,pdfborder={0 0 0}}

% --- Batch metadata ---------------------------------------------------------
% Taken as they stand by the rendering script, which shows them at the head of
% the reading view. They fix what the file claims to cover. The macro names are
% the parent project's, so that the scripts stay shared; \folder here means the
% volume's slug — `fr-9115` — and \pages counts Gallica views.

\newcommand{\folder}[1]{\gdef\grothFolder{#1}}
\newcommand{\batch}[1]{\gdef\grothBatch{#1}}
\newcommand{\pages}[2]{\gdef\grothFirst{#1}\gdef\grothLast{#2}}
\newcommand{\foldertitle}[1]{\gdef\grothFoldertitle{#1}}

% The holder's own shelfmark, printed as they write it: « Français 9115 ».
\newcommand{\shelfmark}[1]{\gdef\grothShelfmark{#1}}

% The dating, copied verbatim from the holder's catalogue — never inferred
% here. « XIXe siècle » is the BnF's; a pass that dates a leaf from its content
% says so in a \note{}, as its own inference.
\newcommand{\dating}[1]{\gdef\grothDating{#1}}

% The status notice, printed in the title block and repeated as a diagonal
% watermark on every page of the PDF. The manuscripts are in the public
% domain; what this line declares is not a right but a quality — an
% unverified first pass by a machine. The skills write the line; a file
% without it gets no watermark, so leaving it out is visible in the output.
\newcommand{\watermark}[1]{\gdef\grothWatermark{#1}}

\gdef\grothFolder{?}\gdef\grothBatch{}\gdef\grothFirst{?}\gdef\grothLast{?}%
\gdef\grothFoldertitle{}\gdef\grothShelfmark{}\gdef\grothDating{}\gdef\grothWatermark{}

% --- The résumé -------------------------------------------------------------
\newenvironment{resume}
  {\small\setlength{\parskip}{0.45em}}
  {\par\medskip\hrule\medskip}

% --- Keywords ---------------------------------------------------------------
% In English, closing the modernised reading's résumé: the modern vocabulary
% under which to search for what the leaves construct. The manifest extracts
% them to tag the volume — this line is the single source of the tags.
\newcommand{\keywords}[1]{\par\smallskip\noindent
  {\footnotesize\textsc{Keywords} --- \textit{#1}}\par}

% --- The critical apparatus -------------------------------------------------

% The Gallica view number, in the margin. This is the anchor that lets a
% disagreement over a formula be settled by looking at *one* image rather than
% a volume of 750. The site uses it to turn the facsimile as the transcription
% is scrolled. A view is one scanned image: usually one side of a leaf.
\newcommand{\page}[1]{%
  \marginnote{\footnotesize\color{gray!70}\textsf{v.\,#1}}%
  \label{page:#1}\ignorespaces}

% The BnF's pencilled foliation, where the leaf carries it and a pass can read
% it: « 198r ». This is what the literature cites — « ff. 198r–208v » — and
% the one line that turns a citation into a view. Recorded, never inferred:
% a folio a pass cannot read is not written.
\newcommand{\folio}[1]{%
  \marginnote{\footnotesize\color{gray!50}\textsf{f.\,#1}}\ignorespaces}

% The modernised reading's coarser counterpart to \page{}: which views a
% section is drawn from.
\newcommand{\pagerange}[2]{%
  \marginnote{\footnotesize\color{gray!70}\textsf{v.\,#1--#2}}%
  \ignorespaces}

% Illegible: never guessed.
\newcommand{\ill}{\textcolor{red!60!black}{[\,\ldots\,]}}

% A doubtful reading: the word is given, and flagged as uncertain.
\newcommand{\uncertain}[1]{\uline{#1}}

% An editorial addition — a missing letter, an implied word.
\newcommand{\add}[1]{\textcolor{blue!50!black}{[#1]}}

% Struck out by the writer. What was crossed out is part of the document.
\newcommand{\struck}[1]{\sout{#1}}

% The transcriber's note, on the state of the page or the mathematics.
\newcommand{\note}[1]{\par\smallskip{\small\color{gray!60!black}\textit{[#1]}}\par\smallskip}

% A marginal note of hers, distinct from the body of the text.
\newcommand{\marginal}[1]{\par\smallskip\noindent\hspace{1em}%
  {\small\color{gray!60!black}#1}\par\smallskip}

% --- Title ------------------------------------------------------------------
% The title block is French, like the documents themselves.

\AddToHook{shipout/background}{%
  \ifx\grothWatermark\empty\else
    \begin{tikzpicture}[remember picture, overlay]
      \node[rotate=52, scale=3.4, align=center, text=gray!18,
            font=\sffamily\bfseries]
        at (current page.center) {\grothWatermark};
    \end{tikzpicture}%
  \fi}

\AtBeginDocument{%
  \begin{center}
    {\large\bfseries Manuscrits de mathématiciens (Gallica) ---
      \ifx\grothShelfmark\empty\grothFolder\else\grothShelfmark\fi}\\[2pt]
    {\itshape\grothFoldertitle}\\[4pt]
    {\small vues \grothFirst--\grothLast
      \ifx\grothBatch\empty\else\ (lot \grothBatch)\fi}\\[2pt]
    \ifx\grothDating\empty\else
      {\small\color{gray!60!black}Datation du catalogue : \grothDating}\\[2pt]
    \fi
    \ifx\grothWatermark\empty\else
      {\small\bfseries\color{red!45!gray}\grothWatermark}\\[2pt]
    \fi
    {\footnotesize\color{gray!60!black}Transcription établie d'après les images IIIF de Gallica.
     Source gallica.bnf.fr / Bibliothèque nationale de France.\\
     Les lectures incertaines sont soulignées, les illisibles notées [\,\ldots\,],
     les ajouts entre crochets ; \textsf{v.} numérote les vues de Gallica,
     \textsf{f.} la foliotation au crayon de la BnF.}
  \end{center}
  \vspace{1em}}

\endinput


\folder{fr-12357}
\shelfmark{Français 12357}
\batch{3}
\pages{41}{56}
\dating{1601-1700}
\watermark{Lecture automatique\\non vérifiée}
\foldertitle{« Remarques tirées du livre de l'Harmonie universelle du P. MERSENNE, ainsy qu'il les avoit escrites, de sa main, à la marge et aux feuillets blancs devant et derrière dudit livre ».}
\begin{document}

\page{41}
\note{Page de gauche (verso) ; aucun chiffre en tête. La première ligne continue la dernière de la vue 40 (« et que d'un »). Les fins de ligne passent dans le pli de la reliure ; les lettres cachées ne sont complétées par \add{} que lorsque le mot est certain.}

tuyau de 12 pieds de hauteur, elle monte par le trou d'vne ligne dans $1''$, elle ne pourra monter aussy viste que la bale d'vne harquebuse si le tuyau n'est ha\add{ut} de 43200 pieds, supposé que la bale face 100 toises dans vne seconde minute, c'\add{est} a dire 60 fois plus viste que l'eau qui monte dix pieds en sortant d'vn tuyau de 12 pieds de haut, car l'autre tuyau doibt auoir sa hauteur en raison doublée de \ill{} a 60 ou de 12 a \struck{720} 720.
\note{Après « doublée de », le premier terme du rapport est caché dans le pli. Le « 720 » biffé est récrit aussitôt, identique.}

l'eau sortant d'vn canal et montant dans vn aut\textsuperscript{re} moindre, par exemple sortant d'\add{un} tuyau d'vn pouce de large et entrant dans vn autre de deux lignes de larg\add{e} cette eau met deux fois autant de temps a y monter co\textsuperscript{e} lors qu'elle monte dans l'air par iect sans canal; car le iect d'vn tuyau de quatre pieds de hau\add{t} dure $1''$ et la conduite par le petit canal $2''$.

p. 6. coroll. 1\textsuperscript{er} ligne 2.

elle seroit perpetuelle dans vn milieu qui n'auroit point de resistence, suyuant le\add{s} raisonnem\textsuperscript{ts}. faits ailleurs.

lig. 13.

ayant repesé la vessie et le plomb, ie trouue qu'il est seulem\textsuperscript{t}. 1890 fois plus pesa\add{nt} qu'elle, partant le plomb, partant le plomb estant descendu de 1890 pieds ou 315 t\ill{}
\note{« partant le plomb » est écrit deux fois par le copiste. La phrase s'arrête au pli sur « 315 t… » (315 toises valent 1890 pieds) ; la ligne suivante passe à la remarque sur la p. 7, sans que la phrase soit achevée sur la page.}

p. 7 lig. 5.

elles ne sont pas proprem\textsuperscript{t} plus legeres car elles sont poussées par la chaleur ou autre agent exterieur, et sans cela elles retombent.

p. 7. coroll. 5

voyez l'experience contre Galilée page 67 de ses nouuelles pensées apres le 15\textsuperscript{me}. arti\add{cle} ou il est monstré que le plomb debvroit tomber de 42 pieds de hauteur dans l'ea\add{u} dans le temps de $2''$, au lieu qu'il ne tombe que de 12 pieds, parce que l'eau qui pese 12 fois moins que le plomb, ne debvroit empescher qu'vn 12\textsuperscript{me} de sa cheu\add{te} et l'air $\frac{1}{480}$ partye de son mouuem\textsuperscript{t}. dans le vuyde.
\note{« page 67 » : le 6 a la forme bouclée que ce copiste donne partout à ce chiffre.}

p. 9. lig. 25.

J'ay faict et mis ceste experience dans les reflexions physicomathematiques \ill{} chap. 22.

Ibid. coroll. 8.

enfin le 4\textsuperscript{e} Juin 1638 j'ay trouué que le plomb est a la vessie, co\textsuperscript{e} 1890 a 1. voyez la 29 et 32. prop. de l'hydraulique pneumatique, ou l'air est exactem\textsuperscript{t} pesé, et \ill{} pese moins que l'eau 1300 fois
\note{Le mot illisible est couvert par une tache d'encre.}

p. 10. seconde obseruation

Il y a encore vn aut\textsuperscript{re} mouuem\textsuperscript{t} composé du naturel et du violent, lequel descrit u\add{ne} ligne parabolique telle que ABC en ceste façon. Que le mobile B, ayt vne vistess\add{e} tous\textsuperscript{rs}. egale en se mouuant vers E, par laquelle il descriue BA, par le moyen du mouuem\textsuperscript{t}. par BD, que sa pesanteur luy faict f\textsuperscript{e}, si l'on suppose le mobile en A estre de mesme vistesse tous\textsuperscript{rs}. egale vers D, et qu'il soit aussy poussé en haut d'vne fa\add{çon} contraire a celle qu'il estoit descendu de B en D, en repassant par tous les degrez d\add{e} vistesse qu'il auoit acquis en descendant, ce qui arriuera si ayant descendu en D, \ill{} remonte par la mesme force qu'il y a acquise, car la pesanteur destruira ces degrez \ill{} vistesse acquise depuis B iusques a D, en mesme proportion qu'ilz ont esté acqu\add{is}. Cecy posé il est euident, que le mobile A, descrira la demye parabole, \struck{et} AB, qu'estant en B, puis qu'il continue de la mesme vistesse a se mouuoir de B en F, et \ill{} les degrez qu'il a acquis en montant, se destruisent au contraire sens en descendant de B \ill{} D, il est certain qu'il \struck{destruira} descrira l'aut\textsuperscript{re} partie de la parabole BC, de sorte que si on la repoussoit en haut co\textsuperscript{e}. de A, et qu'il retournast de mesme vistesse de C vers A, \ill{} repasseroit par la mesme parabole CDA, ou on referoit vne semblable de C, vers la ma\add{in} droite, ou s'il continuoit il alongeroit la parabole vers H. mais tout cecy suppose qu'il n'y \ill{} nul empeschem\textsuperscript{t} de l'air, lequel altere merueilleusem\textsuperscript{t}. ceste proportion.
\note{Les \ill{} de ce paragraphe sont des fins de ligne cachées dans le pli de la reliure (un mot court à chaque fois). Figure dans la marge gauche : une demi-ellipse ou parabole ABC au-dessus d'une base A D C, la tangente horizontale E B F passant par le sommet B, la verticale BD, et la lettre H sous C, où la courbe se prolonge.}

Il faut corriger tout cecy, puisq. les \uncertain{flesches} conuainquent que la descente est plus long t\ill{} a se f\textsuperscript{e} que la montée, co\textsuperscript{e} i'ay desia remarqué es lieux cotez au commencem\textsuperscript{t} de cest\add{e} obseruation. voyez la balistique

\note{Réclame au bas de la page : « p. 13 lig. 25. », répétée en tête de la vue 42.}

\page{42}
\note{Page de droite (recto), chiffrée « 20 » en haut à droite, d'une main plus fine et d'un autre tracé que le texte : un chiffre par feuillet, au recto seulement, qui court jusqu'à 26 à la vue 55 ; cette passe ne le donne pas pour la foliation de la bibliothèque (voir l'en-tête). Au pied de la page, à droite, un « 3 » isolé, d'une autre plume ; sa nature (signature de cahier ?) est une inférence.}

p. 13. lig. 25.

l'on pourroit f\textsuperscript{e} l'essay auec vne arbaleste à flesche, en obseruant le temps qu'elle employe à tomber par son propre poids depuis le haut d'vne tour, et celuy qu'elle employe en tombant par la proiection de l'arc, apres auoir mesuré combien de temps dure la proiection de bas en haut, iusqu'au haut de la tour.

p. 15. obseruaõn 4\textsuperscript{me}. lig. 7.

A cecy se raporte ce que fait le s\textsuperscript{r}. \uncertain{van Etchino}, cousin de m\textsuperscript{r}. Huygens de Zuylichen, lequel siflant dans vn verre en retirant, et no\textsuperscript{n}. en poussant son vent fait f\textsuperscript{e} toute sorte de tons aud\textsuperscript{t}. verre, il fait la mesme chose dans les cloches a ce que l'on m'escrit d'vtrect.
\note{Le nom est lu tel quel, lettre à lettre ; « Zuylichen » aussi, pour Zuylichem.}

p. 16. auant la 5\textsuperscript{me} obseruation

Galilée dit auoir obserué dans sa premiere iournée, que les gouttes de l'eau qui fremit, se fendent en 2 partyes, lors que le son du verre saute a l'octaue en haut

\marginal{n\textsuperscript{l}}
5\textsuperscript{e} obseruation

voyez les lettres de m\textsuperscript{r}. descartes sur ce subiect, ou il determine ce qui en est.

p. 17. au commencem\textsuperscript{t}.

\marginal{n\textsuperscript{a}}
L'on a monstré les erreurs de la \uncertain{distraõn} de m\textsuperscript{r}. de beaugrand sur ce suiect, et ce que l'on en peut dire au fonds: voyez le liure de m\textsuperscript{r} de la brosse, et les lettres de m\textsuperscript{r}. descartes sur ce subiect. voicy la refutation
\note{Les deux signes de la marge gauche, « n » suivi d'un l ou d'un a suscrit, sont reproduits tels quels ; leur sens (nota ?) n'est pas établi. « distraõn » : un mot abrégé, d, i, s, t, r, a, o, n avec un trait suscrit ; il n'est pas développé ici.}

Soit la figure qui est a la marge dans laquelle A, represente le centre de la terre, G, celuy d'vne balance; dont FD sont les deux bras, puis mettant vn poids au point F, et vn aut\textsuperscript{re} attaché au point D, qui pend plus bas iusqu'au point E, ou le poids pese d'autant moins qu'il est plus proche du centre.
\note{Figure dans la marge gauche : un triangle F D A, le fléau FD horizontal, G sur FD, la verticale GA passant par H, et la droite FE, E sur DA, qui coupe GA en H.}

Contre cecy, l'on dit 1\textsuperscript{o}. qu'encore que ce poids ainsy posé pesast moins au regard des autres poids, qui luy seroient opposez dans cette balance, il ne s'ensuit pas pō cela qu'il doiue moins peser estant consideré tout seul hors de la balance.

2\textsuperscript{o} le principe d'Archimede et de Pappus, a scauoir que le centre de deux corps pesans ioints ensemble diuise la ligne droite qui conioint leurs centres, en raison reciproque de leurs pesanteurs, n'est vray, ny entendu par ces autheurs, qu'en cas que les corps pesans tendent en bas par lignes paralleles, et sans s'incliner vers vn mesme point, et partant l'autheur de la geostatique qui suppose le contraire ne prouue rien.

3\textsuperscript{o} si sa proposition estoit vraye ce qu'il dit du centre de grauité seroit faux: par exemple si le poids F, et D, sont egaux, leur commun centre, selon archimede sera en G, au lieu que, selon l'autheur, quand le poids D, est plus proche du centre de la terre que le poids F, ce centre doibt estre entre F et G, et quand il est plus esloigné, ce centre doibt estre entre G et D.

4\textsuperscript{o} ayant supposé le poids 1\textsuperscript{o}. estre au poids \uncertain{B}\textsuperscript{o}, lors qu'ilz sont en pareille distance du centre de la terre, co\textsuperscript{e} la ligne EH, est a FH, il ne les met pas a pareille distance, mais l'vn au point F, et l'autre au point E, et puis il suppose que le point H, est le centre de grauité, de mesme que s'ilz estoient a egale distance, et ainsy pō prouuer que ce changem\textsuperscript{t} de distance, change la pesanteur, il suppose qu'il ne la change point.
\note{« le poids 1\textsuperscript{o}. estre au poids B\textsuperscript{o} » : lu tel quel ; la figure n'a pas de B, et la lettre pourrait être un 8. Une tache d'encre touche le signe.}

5\textsuperscript{o} il s'apuye sur ce que F, est plus esloigné d'A, que n'est E, de sorte que si on le suppose plus proche, et que l'on recoyue tout le reste de son discours comme vray, l'on en conclurra tout le contraire de ce qu'il conclud, et toutesfois en construisant sa figure il donne la liberté d'y f\textsuperscript{e} la ligne AF, de telle grandeur qu'on voudra

voyez la 6\textsuperscript{me}. raison p. 61 de l'vtilité de l'harmonie prop. 18.

p. 18. en la 6\textsuperscript{e} obseruation

voyez les figures de ces instruments attachés a la 6\textsuperscript{me} proposition du 5\textsuperscript{me} liure des instruments, auec les lieux de leurs trous et leurs nombres

\page{43}
\note{Page de gauche (verso) ; aucun chiffre en tête. La dernière colonne du tableau passe dans le pli de la reliure : son titre et le dernier chiffre de plusieurs de ses nombres y sont cachés.}

p. 22. obseruation 9\textsuperscript{me}

\[
\begin{array}{l|ll|ll|ll|ll|l}
\text{Accords du monocorde} & \text{parfaits} & \text{A} & \text{Aprochants} & \text{B} & \text{Imparfaits} & \text{C} & \text{Du tout esloignez} & \text{D} & \text{Accord de Ga…} \\
\hline
\text{Denominaõn ou diuision du monocorde en} & & 100000 & & 100000 & & 100000 & & 100000 & 100000 \\
\hline
\text{Tertia minor comprend les} & \frac{5}{6}\ \text{ou} & 83333\frac{1}{3} & \frac{1024}{1215}\ \text{ou} & 84280 & \frac{27}{32}\ \text{ou} & 84375 & \frac{64}{75}\ \text{ou} & 85333 & 84090 \\
\quad\text{Concordent} & \text{G.}\ \natural\text{.\ C} & \text{E. Fx} & & & & & & & \\
\quad\text{Discordent} & \text{---} & \text{---} & \text{F} & \text{---} & \text{C gx. a} & \text{Cx. d.} & \text{B} & \text{dx.} & \\
\text{Difference auec la parfaite} & & \text{concordance} & +\ \text{basse} & -\ 1124 & +\ \text{basse} & -\ 1235 & +\ \text{basse} & 2344 & +\ \text{basse}\ \ 900 \\
\hline
\text{Tertia maior comprend les} & \frac{4}{5}\ \text{ou} & 80000\ \text{ou} & \frac{405}{512} & 79102 & \frac{64}{81}\ \text{ou} & 79012 & \frac{25}{32}\ \text{ou} & 78125 & 79370 \\
 & \text{G.A.B.C} & \text{D dx} & \text{Gx} & \text{Cx} & \text{---} & \text{E F} & & \natural\ \ \text{Fx} & \\
\text{Difference d'auec la parfaite} & & & +\ \text{haut} & 1124 & \text{haut} & 1235 & \text{haut} & 2344 & \text{haut}\ \ 789 \\
\hline
\text{Quarta comprend les} & \frac{3}{4}\ \text{ou} & 75000 & \frac{512}{675}\ \text{ou} & 75852 & \frac{20}{27}\ \text{ou} & 74074 & & & 74915 \\
\quad\text{Concordent G. Gx.} & \text{A.B.}\natural & \text{C Cx. D, Fx.} & & & & & & & \\
\quad\text{Discordent} & & & \text{Dx} & & \text{E F.} & & & & \\
\quad\text{Difference} & & & \text{basse} & 1123 & \text{haulte} & 1235 & & & \text{haulte}\ \ 113 \\
\hline
\text{Quinta comprend les} & \frac{2}{3}\ \text{ou} & 66667 & \frac{675}{1024}\ \text{ou} & 65918 & \frac{27}{40}\ \text{ou} & 67500 & & & 6674… \\
\quad\text{Concordent G.}\natural\text{.C.} & \text{Cx.D.Dx.E.F Fx.} & & & & & & & & \\
\quad\text{Discordent en} & & & \text{Gx} & & & \text{A.B.} & & & \\
\quad\text{Differences} & & & \text{haulte} & 1123 & \text{basse} & 1235 & & & 112 \\
\hline
\text{Sexta minor comprend} & \frac{5}{8}\ \text{ou} & 62500 & \frac{256}{405}\ \text{ou} & 63210 & \frac{81}{128}\ \text{ou} & 63282 & \frac{16}{25}\ \text{ou} & 64000 & 62996 \\
\quad\text{Concordent en} & \text{G.}\ \natural & \text{Cx.D.E.Fx} & & & & & & & \\
\quad\text{Discordent en} & & & \text{C\ \ F} & & \text{Gx} & \text{A} & \text{B.} & \text{Dx} & \\
\quad\text{la difference} & & & \text{basse} & 1123 & \text{basse} & 1235 & \text{basse} & 2344 & \text{basse}\ \ 789 \\
\hline
\text{Sexta minor comprend} & \frac{3}{5}\ \text{ou} & 60000 & \frac{1215}{2048}\ \text{ou} & 59326 & \frac{16}{27}\ \text{ou} & 59259 & \frac{75}{128}\ \text{ou} & 58994 & 59460 \\
\quad\text{Concordent en G. A} & \text{B.D.E} & \text{Dx} & & & & & & & \\
\quad\text{Discordent en} & & & \text{Gx} & & \natural & \text{C.E.F.} & & \text{Cx.\ Fx} & \\
\text{Differences des concordances auec la parfaite} & & & \text{haulte} & 1123 & \text{haulte} & 1235 & \text{haulte} & 2344 & \text{haulte}\ \ 900
\end{array}
\]
\note{Le tableau est reproduit case par case. Les en-têtes « parfaits », « Aprochants », « Imparfaits », « Du tout esloignez » et « Accord de Ga… » sont écrits sous des arcs qui coiffent chacun deux colonnes (la fraction et sa valeur sur une corde de 100000) ; la dernière colonne, dont le titre se perd dans le pli, n'en a qu'une. Ses derniers chiffres, cachés dans le pli sur cette vue, se lisent au bord gauche de la vue 44 (100000, 62996, 59460) ; le cinquième chiffre de 6674… n'y paraît qu'à moitié, peut-être un 2. Dans la première ligne, « Acco » est biffé devant « Denominaõn ». Une lettre suivie d'un trait en forme de x (Cx, Dx, Fx, Gx) est transcrite avec un x : c'est, par inférence, la diésis de la note ; le signe en forme de 4 est transcrit par le bécarre, comme dans la vue 21 et suivantes. Lectures douteuses : « E. Fx » (tierce mineure, colonne A), où la seconde lettre pourrait être un f ou un s bouclé ; « C gx. a » (tierce mineure, colonne C), où le C est biffé et la lettre suivante est lue g suivie du même trait ; « B.D.E » (dernière sexte), dont la troisième lettre, bouclée, pourrait être un A. Dans la tierce majeure, le « ou » de la colonne B est écrit à la suite de 80000. La dernière ligne, de rapport $\frac{3}{5}$, est intitulée comme la précédente « Sexta minor » : c'est la sexte majeure, et le titre est reproduit tel qu'il est écrit.}

Les nombres des colonnes A,B,C,D, demonstrent les accords selon le monocorde de Ptolomee, Boetius, Zarlin et caet. L\ill{} Colonne A, denote les concordances parfaites, les autres sont imparfaites et la derniere colonne est l'accord de Ga\ill{} par moyennes proportionelles.

Les 12 semitons ont leurs octaues accordées parfaitem\textsuperscript{t} lesquelles auec les 40 accords parfaits qui se trouuent en la 1\textsuperscript{re} colonne, ne font que 29 concordances parfaites, si ce n'est qu'on veuille aussy passer 8 aprochantes pour consonances \struck{par} parfaites cela estant admis et receu faudra aussy admettre les concordances de Gallé, puisqu'elles aprochent encore beaucoup plus pres des parfaites que lesd. aprochan\add{tes} car ses quintes et quartes ne different des vrayes et tres parfaites que de $\frac{113}{100000}$ et les tierces maieures et s\add{extes} mineures de $\frac{789}{100000}$, les tierces mineures et sextes maieures de $\frac{900}{100000}$ au lieu que les aprochantes different de $\frac{1123}{100000}$ notez que la tierce maieure ne differe pas tant de l'accord parfait que la tierce mineure et cepend\add{ant} les facteurs d'orgues accoustumés aux erreurs de \ill{} la iugent au sentim\textsuperscript{t} de l'ouye plus imparfaite, au lieu de ceux qui n'exercent que la voix, ou qui iouent d'aut\textsuperscript{res} instrum\textsuperscript{ts} co\textsuperscript{e} luth, theorbe, viole de gamb\add{e} cornette et caet: la iugent du contraire par le mesme sentiment, car leur voix les porte naturellem\textsuperscript{t} pa\add{r} tons et semitons egaux, voire leur luth font de mesme que leur voix excedant tant soit peu en la tierce maieure, et s'aprochant de si pres en la tierce mineure, co\textsuperscript{e} est l'excez de la maieure, soit pō \uncertain{vng} temp\ill{} mechanique diuisée l'octaue selon le commun en 53 comma desquels 14 font la tierce mineure et 17 fon\add{t} la maieure, si les tons fussent diuisez egalem\textsuperscript{t}, chascun contiendroit $8\frac{5}{6}$ de comma et chasque demyton 4\ill{} de comma par ainsy la tierce mineure debvroit contenir $13\frac{1}{4}$ de comma, et la tierce maieure $17\frac{2}{3}$ laquelle \ill{} rehaussée que de $\frac{2}{3}$ d'vn comma au dessus de sa parfaite et la tierce mineure est afayblie et abaissée de \ill{} d'vn comma au dessoubs de sa parfaite qui est $\frac{1}{12}$ qu'elle a de plus, pō difference de cecy, on cognoist q\ill{}
\note{Une tache d'encre couvre la fin du « par » biffé. « Gallé » : lu tel quel, avec un trait sur l'e ; c'est le nom qu'abrège l'en-tête de la dernière colonne. Après « aux erreurs de », un mot court lu « lead » ou « lrab », non résolu. « soit pō vng temp… » : la fin de la ligne passe dans le pli. Les autres \ill{} de ce paragraphe sont des fins de ligne cachées dans le pli : la fraction qui suit « chasque demyton 4 », un mot court avant « rehaussée », la fraction qui suit « abaissée de ». Le dénominateur de $\frac{1123}{100000}$ est suivi d'une barre oblique.}

\note{Réclame au bas de la page : « que le Ingenu », reprise en tête de la vue 44.}

\page{44}
\note{Page de droite (recto), chiffrée « 21 » en haut à droite, de la même main fine que le « 20 » de la vue 42. La première ligne continue la réclame de la vue 43 (« que le Iugem… »). Au bord gauche de la vue paraît le bord droit de la vue 43, avec la dernière colonne de son tableau (voir la note de cette vue). Le quadrillage pâle qui traverse la moitié inférieure de la page est le carré magique de la vue 45, vu par transparence.}

que le iugem\textsuperscript{t} de l'ouye des facteurs d'orgues est erroné et \uncertain{peruerty} par accoustumance et continuation d'entendre tous\textsuperscript{rs} la mesme erreur, l'ayant admis co\textsuperscript{e} chose necessaire laquelle leur est en fin passée en nature, voyla pourquoy ce n'est pas merueille s'ilz iugent la tierce mineure des tons egaux (qui a esté demonstrée moins parfaite que la maieure) meilleure que lad\textsuperscript{t}. maieure au contraire ceux qui ne possedent que la musique vocale qui iugent plus sainem\textsuperscript{t} \uncertain{confessant} tous la maieure plus agreable que la mineure.
\note{« confessant » : les lettres se lisent « coffissant », sans signe d'abréviation visible.}

p. 25. au commencem\textsuperscript{t}

L'on a semblablem\textsuperscript{t} trouué la proportion de ceste figure de roulette, quelque raison qu'il y ayt entre la circonference de la roulette, et le chemin droit qu'elle fait, dont ie mets icy quelques exemples, et l'on a trouué les tangentes de lad\textsuperscript{t}. figure, voyez les lettres de m\textsuperscript{r}. descartes.

L'on a encore trouué le raport desd\textsuperscript{t}. lignes auec les figures ellyptiques et toutes autres roulantes sur vn plan horizontal, et la raison du solide de la roulette à vn cylindre de mesme hauteur, lors qu'elle tourne sur vne ordonnée, car elle est comme 8 à 5, voyez les demonstrations de la roulette, et de ses tangentes dans les lettres de m\textsuperscript{r}. descartes. voyez aussy, la preface des hydrauliques et les centres de grauité de la roulette ont aussy esté trouuez et de ces solides et cat.
\note{« encore » porte un petit trait suscrit à la fin, peut-être un s. « dont ie mets icy quelques exemples » : aucun exemple ne suit sur la page.}

p. 26. au commencem\textsuperscript{t}.

Si sumatur binarius, vel quilibet alius numerus ex binarij multiplicatione productus, modo sit talis vt si tollatur vnitas ab eius triplo fiat numerus primus, Item si tollatur vnitas ab eius sextuplo fiat numerus primus, et deniq; si tollatur vnitas ab eius quadratj octodecuplo fiat numerus primus, ducaturq; hic \struck{vlt} vltimus primus per duplum numerj assumptj, fiet numerus, cuius partes aliquotae dabunt aliū numerū, qui vice versâ partes \struck{aliq} aliquotas habeat aequales numero praecedentj. Sic assumendo tres numeros, 2, 8 et 64 habeo, haec tria paria numerorū, aliaq; infinita inueniri possunt parj modo ex Cartesio

284: 220

18416: 17296

9437056: 9363584
\note{Le copiste écrit la terminaison -ur (« sumatur », « tollatur ») par un signe en forme de z, développé ici sans autre marque. Le premier chiffre de 18416 et de 9437056 touche la marge ; le 9 de 9437056 est souligné.}

xiij obseruation. lig. 4.

cette methode ne vaut rien que pō \uncertain{vn} deux nombres ausquels elle a esté accommodée sans que par elle on en puisse trouuer d'autres.

Il n'y en a dans tous les nombres, que les 6 qui suyuent, qui ayent leur p\textsuperscript{ties} aliquotes doubles de leur nombres.

120

672

523776 \quad Autres nombres dont les partyes aliquotes sont le double

51001180160

1476304896 \quad trouué par m\textsuperscript{r}. descartes

459818240 \quad multiplié par 3

1379454720 \quad dont les partyes aliquotes sont le triple
\note{Une accolade réunit 120, 672 et 523776 et renvoie à « Autres nombres dont les partyes aliquotes sont le double », écrit à droite entre deux accolades. 523776 est écrit « 5.23776 ». Un petit 3 est écrit au-dessus des deux derniers chiffres de 1379454720, qui est le produit de la ligne précédente par 3.}

les \struck{huit}\textsuperscript{9} qui suyuent ont aussy des partyes aliquotes qui sont le triple

30240 / 32760 / 23569920 / 45532800 / 142990848 / 1379454720 / 43861478400 / 66433720320 / 4\uncertain{0}3031236608
\note{Chacun de ces neuf nombres porte au-dessus de lui un petit chiffre d'ordre, de 1 à 9. Le deuxième chiffre du dernier nombre, sur lequel tombe le chiffre d'ordre 9, peut se lire 0 ou 8.}

nombres dont les p\textsuperscript{ties} aliquotes sont le quadruple / 14182439040 / 508666803200 / 30823866178560 /

Quand aux nombres qui se refont mutuellem\textsuperscript{t} voyez la reigle de m\textsuperscript{r}. descartes pō les trouuer en voicy trois
\[
\begin{array}{|lcl|}
284 & — & 220 \\
18416 & — & 17296 \\
9437056 & — & 9363584
\end{array}
\]

voyez la preface sur les nouuelles pensées de Galilée

Les nombres suyuans 2147483648, 13, 17, 31, 41, 43, 61, 83, 163, 257, 307, 331, 467, 613, 1093, 1331, 2401, 2801, 65537, 299209, 594323. multipliez l'vn par l'autre, donnent vn quadruple, lequel multiplié par 5, donne vn nombre, dont les p\textsuperscript{ties} sont le quintuple; mais il n'est pas diuisible par 5, et si l'on multiplie vn sextuple par 7, non diuisible par 7, il produira vn septuple.
\note{Ce paragraphe est écrit dans un cadre, dans la moitié droite de la page, en face des listes de nombres.}

\page{45}
\note{Page de gauche (verso) ; aucun chiffre en tête. Les nombres des carrés portent presque tous un point sous un de leurs chiffres, marque du copiste dont le sens n'est pas établi ; ces points ne sont pas reproduits. À droite, au-delà du pli, paraissent les premières cases d'une petite table de la vue 46 (« 1. », « 8. », « 12. », « 13. », « 16: »), qui ne sont pas de cette page.}

Au papier blanc apres tous les liures de l'harmonie

Quarré de 22 contenant ceux de 16, de 12, et de 10 celuy qui fait 12 a 241, 242, 243 et 244 en ses 4 coins ceux de 10 sont 33, 46, 452, et 439: et ceux de 16 sont 55, 251, 234 et 433
\note{« 243 » : le dernier chiffre est à demi couvert par le « et » qui suit.}

\[
\begin{array}{|c|c|c|c|c|c|c|c|c|c|c|c|c|c|c|c|c|c|c|c|c|c|}
\hline
76 & 145 & 304 & 166 & 355 & 144 & 324 & 164 & 149 & 346 & 349 & 150 & 400 & 143 & 352 & 351 & 153 & 362 & 110 & 344 & 160 & 288 \\
\hline
373 & 407 & 406 & 337 & 89 & 94 & 87 & 330 & 300 & 187 & 93 & 372 & 383 & 347 & 96 & 380 & 106 & 82 & 158 & 189 & 192 & 327 \\
\hline
334 & 151 & 80 & 148 & 123 & 333 & 343 & 109 & 122 & 322 & 382 & 113 & 358 & 138 & 315 & 128 & 332 & 348 & 378 & 293 & 381 & 104 \\
\hline
184 & 395 & 359 & 55 & 276 & 77 & 224 & 418 & 429 & 412 & 259 & 432 & 58 & 121 & 124 & 272 & 74 & 415 & 234 & 114 & 102 & 301 \\
\hline
125 & 292 & 388 & 286 & 52 & 249 & 157 & 428 & 56 & 262 & 435 & 53 & 427 & 254 & 257 & 265 & 252 & 248 & 199 & 316 & 169 & 165 \\
\hline
177 & 306 & 190 & 289 & 215 & 241 & 450 & 449 & 448 & 40 & 41 & 42 & 219 & 247 & 245 & 246 & 242 & 256 & 210 & 367 & 297 & 118 \\
\hline
183 & 354 & 313 & 419 & 422 & 191 & 33 & 390 & 277 & 281 & 278 & 284 & 269 & 282 & 285 & 46 & 294 & 63 & 66 & 182 & 84 & 339 \\
\hline
323 & 309 & 384 & 258 & 214 & 283 & 446 & 5 & 1 & 9 & 483 & 482 & 479 & 478 & 8 & 34 & 202 & 434 & 64 & 101 & 176 & 162 \\
\hline
140 & 86 & 83 & 69 & 417 & 217 & 60 & 16 & 18 & 475 & 464 & 473 & 11 & 14 & 468 & 426 & 268 & 71 & 413 & 399 & 402 & 345 \\
\hline
320 & 97 & 329 & 421 & 51 & 279 & 39 & 13 & 470 & 29 & 465 & 19 & 469 & 463 & 17 & 441 & 206 & 271 & 227 & 193 & 156 & 360 \\
\hline
127 & 120 & 377 & 75 & 230 & 436 & 45 & 32 & 27 & 457 & 466 & 4 & 30 & 459 & 454 & 451 & 49 & 255 & 410 & 371 & 85 & 375 \\
\hline
404 & 100 & 81 & 312 & 235 & 47 & 59 & 460 & 455 & 26 & 28 & 458 & 15 & 31 & 453 & 440 & 438 & 250 & 173 & 305 & 385 & 180 \\
\hline
366 & 181 & 79 & 264 & 232 & 280 & 424 & 461 & 23 & 456 & 22 & 20 & 472 & 462 & 24 & 61 & 205 & 263 & 211 & 393 & 317 & 119 \\
\hline
146 & 401 & 303 & 274 & 222 & 218 & 425 & 476 & 471 & 10 & 12 & 467 & 474 & 21 & 25 & 44 & 267 & 253 & 221 & 172 & 131 & 302 \\
\hline
307 & 174 & 394 & 72 & 414 & 38 & 442 & 477 & 480 & 481 & 3 & 2 & 6 & 7 & 484 & 43 & 447 & 68 & 416 & 111 & 340 & 129 \\
\hline
295 & 188 & 318 & 275 & 229 & 437 & 452 & 95 & 203 & 201 & 204 & 216 & 200 & 208 & 207 & 439 & 48 & 270 & 196 & 167 & 179 & 308 \\
\hline
299 & 115 & 186 & 287 & 225 & 243 & 35 & 36 & 37 & 445 & 444 & 443 & 266 & 238 & 240 & 239 & 244 & 260 & 198 & 386 & 310 & 159 \\
\hline
356 & 325 & 194 & 273 & 237 & 233 & 228 & 67 & 65 & 223 & 50 & 431 & 423 & 231 & 328 & 213 & 236 & 433 & 212 & 91 & 311 & 178 \\
\hline
107 & 175 & 112 & 251 & 209 & 411 & 361 & 57 & 420 & 73 & 226 & 54 & 62 & 364 & 261 & 220 & 408 & 70 & 430 & 326 & 370 & 365 \\
\hline
397 & 350 & 195 & 353 & 147 & 363 & 170 & 155 & 152 & 163 & 392 & 369 & 126 & 331 & 142 & 357 & 368 & 396 & 108 & 78 & 88 & 135 \\
\hline
99 & 290 & 92 & 116 & 403 & 185 & 398 & 376 & 391 & 296 & 105 & 132 & 98 & 379 & 389 & 103 & \ast & 137 & 314 & 298 & 409 & 171 \\
\hline
197 & 374 & 168 & 335 & 338 & 336 & 133 & 321 & 341 & 141 & 134 & 319 & 90 & 117 & 161 & 136 & 342 & 130 & 387 & 139 & 291 & 405 \\
\hline
\end{array}
\]
\note{Le carré de 22 est reproduit case par case, rangée par rangée. Chaque rangée et chaque colonne de la lecture donnée ici fait 5335, et les deux diagonales aussi, à trois cases près dont voici le détail. Rangée 5, colonne 14 : un premier chiffre repassé et taché, lu 2 (254) ; 3 donnerait 354, qui est déjà à la rangée 7, et la somme demande 254. Rangée 6, colonne 18 : le chiffre du milieu, bouclé, peut se lire 5 ou 9 ; 296 est déjà à la rangée 21, et la somme demande 256. Rangée 21, colonne 17 : la case est couverte d'une tache d'encre où l'on devine « …54 » ; le nombre 154, que la somme demande, est écrit en clair en face de cette rangée, de l'autre côté du pli (voir plus bas), et la case est marquée ici $\ast$. Rangées 18 et 19 : les cases 433 (rangée 18, colonne 18) et 430 (rangée 19, colonne 19) ont chacune leur dernier chiffre repassé à l'encre épaisse et sont reliées par deux traits obliques partant de leurs coins ; telles qu'elles sont écrites, la rangée 18 fait 5338 et la rangée 19 fait 5332, et l'échange des deux nombres rétablit les deux sommes, comme le demande aussi l'en-tête, qui donne 433 pour un coin du carré de 16 (rangée 19, colonne 19). Les cases 471 (rangée 14, colonne 9) et 477 (rangée 15, colonne 8) sont tachées ; ces deux nombres sont récrits dans un cadre dans la marge droite, en face de ces rangées. Le dernier chiffre de 212 (rangée 18, colonne 19) est repassé. Sous la dernière rangée sont tracés de petits signes, sous les colonnes 8 (×), 9 (+), 10 (un trait), 12 (+), 17 (×) et 21 (un trait) ; leur sens n'est pas établi.}

\marginal{chasque rang de 22 contient 5335}

\marginal{chasque rang de 16 contient 2056 \struck{fa\ill{}} selon le pere \uncertain{mersenne} mais il est fau\add{x} car il contient 3880}

\marginal{chasque rang de 12 contient 870 selon le pere mersen\add{ne} mais il est fau\add{x} car il contient 2910}

\marginal{chasque rang de \ill{} contient 505 selon le pere mersenne m\ill{} il est \ill{} car il con\add{tient} 2425}

\marginal{471 \quad 477}

\marginal{484}
\note{Ces notes sont écrites dans la marge droite, séparées par des traits, en face du carré ; leurs fins de ligne passent dans le pli. Le chiffre du rang de la quatrième (10, d'après la somme 2425) est caché. « 484 » est encadré, en face de la rangée 21. À sa droite, au-delà du pli, donc sur la page d'en face (vue 46), un « 154 » : c'est le nombre que la tache couvre dans cette rangée. Ces notes corrigent les sommes « selon le pere mersenne » : ce sont les seules lignes de cette batch où le copiste parle de Mersenne à la troisième personne et le reprend.}

Quarré de 9 pō 4 planetes, le 9 pō \add{Lune} le 7 pō \add{Vénus}, le 5 pō \add{Mars} et le 3 pō \add{Saturne}

\[
\begin{array}{|c|c|c|c|c|c|c|c|c|}
\hline
73 & 13 & 53 & 26 & 71 & 24 & 49 & 15 & 45 \\
\hline
57 & 22 & 80 & 8 & 76 & 16 & 78 & 7 & 25 \\
\hline
47 & 72 & 44 & 30 & 40 & 36 & 55 & 10 & 35 \\
\hline
23 & 20 & 54 & 39 & 81 & 3 & 28 & 62 & 59 \\
\hline
17 & 12 & 32 & 5 & 41 & 77 & 50 & 70 & 65 \\
\hline
21 & 18 & 48 & 79 & 1 & 43 & 34 & 64 & 61 \\
\hline
31 & 68 & 27 & 52 & 42 & 46 & 38 & 14 & 51 \\
\hline
63 & 75 & 2 & 74 & 6 & 66 & 4 & 60 & 19 \\
\hline
37 & 69 & 29 & 56 & 11 & 58 & 33 & 67 & 9 \\
\hline
\end{array}
\]

369 : 81

Quarré de 8 pō trois planetes, le 8 pō \add{Mercure} le 6 pō $\odot$ et le 4 pō \add{Jupiter}

\[
\begin{array}{|c|c|c|c|c|c|c|c|}
\hline
31 & 9 & 55 & 12 & 54 & 14 & 52 & 33 \\
\hline
43 & 62 & 5 & 29 & 1 & 35 & 63 & 22 \\
\hline
45 & 27 & 26 & 40 & 41 & 23 & 38 & 20 \\
\hline
50 & 61 & 46 & 18 & 17 & 49 & 4 & 15 \\
\hline
8 & 6 & 16 & 48 & 47 & 19 & 59 & 57 \\
\hline
44 & 37 & 42 & 24 & 25 & 39 & 28 & 21 \\
\hline
7 & 2 & 60 & 36 & 64 & 30 & 3 & 58 \\
\hline
32 & 56 & 10 & 53 & 11 & 51 & 13 & 34 \\
\hline
\end{array}
\]

260 : 64
\note{Chaque rangée, chaque colonne et chaque diagonale du carré de 9 fait 369, et du carré de 8, 260, dans la lecture donnée ici ; « 369 : 81 » et « 260 : 64 », sous chacun, donnent cette somme et le nombre des cases. Les signes des planètes sont ceux de la page.}

\[
\begin{array}{|c|c|c|c|}
\hline
5 & 11 & 10 & 8 \\
\hline
16 & 2 & 3 & 13 \\
\hline
4 & 14 & 15 & 1 \\
\hline
9 & 7 & 6 & 12 \\
\hline
\end{array}
\]

16 : 34

on peut varier le quarré magique de 4 en 880 façons, mais celuy de 3 ne se peut f\textsuperscript{e} que d'vne façon, les no\add{mbres} rompus ne peuuent rien faire à \ill{} parce qu'il faut qu'il y ayt tou\ill{} entre deux mesme difference, je n\ill{} \struck{premier} conte point les 8 \struck{differences} transpositions qui se peuuent faire à chasque \ill{}
\note{Ce paragraphe est écrit à droite du carré de 8, sous le carré de 4 ; les fins de ligne passent dans le pli.}

Bienque aucun nombre des enceintes du quarré de 9 ne passast de l'vn en l'aut\textsuperscript{re}, et que les 4\ill{} de la 1\textsuperscript{re} enceinte demeurassent immobiles, et les 8 nombres compris entre iceux y demeurassent aussy, et changeant seulem\textsuperscript{t} de pl\ill{} et quant aux trois aut\textsuperscript{res} enceintes les nombres des coins demeurant \struck{\uncertain{demesmes}} les mesmes auec le seul changem\textsuperscript{t} de 4 places, et que les n\add{ombres} compris entre les coins demeurent tous\textsuperscript{rs} les mesmes l'on peut changer ceste figure en 6742112993280000 façons, sa\ill{} conter les 8 changem\textsuperscript{ts}. de chasque façon qui se fait en tournant sans dessus dessoubs chasque mesme figure. Q\ill{} à l'absolu changem\textsuperscript{t}. il faut prendre la combination ordre de 81.
\note{« demeurant » est surmonté de deux lettres, lues « ts » ; le mot biffé qui suit n'est pas sûr. Le nombre est écrit « 674211299328 0000 », avec un blanc avant les quatre derniers zéros ; il est reproduit sans ce blanc.}

\note{Réclame au bas de la page, à droite : « voyez le quarré ».}

\page{46}
\note{Page de droite (recto), chiffrée « 22 » en haut à droite, de la même main fine que « 20 » et « 21 ». Le texte continue la réclame de la vue 45 (« voyez le quarré »). Au bord gauche de la vue paraissent les notes marginales de la vue 45 (« chasque rang… »), qui ne sont pas de cette page.}

voyez le quarré de 14 (qui a 196 nombres et dont l'addition des nombres de chasque rang se monte à 1379,) \struck{ostant} en vne lettre de m\textsuperscript{r}. fermat \struck{du} au pere mersenne du p\textsuperscript{r} \ill{} Apvril 1640, Et ostant d'iceluy deux enceintes le quarré de 10 qui reste a 985 en chasque rang, et ostant encores deux enceintes pō auoir le quarré de 6. chascun de ses rangs contiendra 591
\note{« du p\textsuperscript{r} » : lu comme l'abréviation de « premier » ; ce qui suit, au bord du pli, est un signe souligné qui n'est pas lu. La date se lit donc, par inférence, « du premier d'Apvril 1640 ».}

\[
\begin{array}{|c|c|c|c|}
\hline
1 & 14 & 15 & 4 \\
\hline
8 & 11 & 10 & 5 \\
\hline
12 & 7 & 6 & 9 \\
\hline
13 & 2 & 3 & 16 \\
\hline
\end{array}
\]

16 : 34
\note{Ce carré de 4 est tracé dans la marge gauche, en face du paragraphe précédent et du suivant.}

Chasque rang du quarré de 4 contient 34, et ce quarré demeurant tous\textsuperscript{rs}. magique se varie en 16 façons differentes, chascune desquelles se varie encores en 8, de sorte que sa varieté entiere \uncertain{est} de 128, mais toutes ses dispositions differentes sont la combinaison ordre de 16 ascauoir 20922789888000, suyuant la table de la 108 page du liure de la voix
\note{« est » : un o surmonté d'un trait, que le copiste emploie aussi plus bas (« magique est 880 ») là où le sens demande « est ».}

La varieté du quarré de 14 ostant chasque enceinte iusqu'à 4, et le reste demeurant tous\textsuperscript{rs} de mesme contient tout ce nombre 3371056466400\uncertain{0} Et le nombre des changem\textsuperscript{ts} du quarré de 22, d'où l'on puisse oster la 1\textsuperscript{re} enceinte, le reste demeurant esgal est

1332756462138013445206441942056960 multiplié par

1188496086155689918464000\uncertain{0}00000 et le produit par

11411553620219513137885347840000
\note{« 3371056466400\uncertain{0} » : un blanc sépare 33710 du reste ; le dernier chiffre est couvert d'une tache. Dans le premier des trois grands nombres, un petit 0 est écrit au-dessus, entre le 4 et le 5 de « …3445… » ; il n'est pas inséré ici. Dans le deuxième, le quatrième des neuf zéros finaux a une queue qui le fait ressembler à un 6. Les trois nombres sont écrits l'un sous l'autre, en tête de ligne.}

La varieté absolüe de 4, c'est à dire de 16 demeurant tous\textsuperscript{rs}. magique \uncertain{est} 880, et y demeurant aux 4 coins 4 nombres qui se suyuent immediatem\textsuperscript{t}, il y en a 16 seulem\textsuperscript{t}. co\textsuperscript{e}. est 4, 5, 6, 7 chascune de ces 880 façons se peut encores voir en 8 manieres differentes en le tournant sans dessus dessous des costes conuerse et inuerse

Arithmeticis proponitur numerus 2027651281, dicant intra 4 dies num sit numerus primus, aut compositus, quod si facere nequeant, nos intra horae spatium certis regulis, et catholicis dabimus.
\note{« num » : lu tel quel.}

\marginal{1}
vn nombre estant donné, par exemple 108, trouuer combien il a de partyes aliquotes, il faut prendre deux nombres qui facent en se multipliant 108, qui soyent ou premiers, ou puissances par exemple 4, et 27, qui sont vn quarré et vn cube, et prendre leurs nombres analogues ou opposants des dignitez 2 et 3 et adiouster à chacun d'eux l'vnité pō auoir 3 et 4, qui multipliez font 12, dont vn osté nō auons 11, pō les partyes aliquotes
\note{« nō » : « nous », abrégé.}

ou bien pō vne deuxiesme methode sans rien adiouster ny oster de 2 et 3, il faut dire 2 fois 3 font 6, c'est a dire prendre leur pareilles ou a vn rectangle, et puis les adiouster on aura 5, or 5 et 6 font 11 partyes aliquotes co\textsuperscript{e} deuant

\marginal{2}
le nombre des partyes aliquotes estant donné, trouuer le moindre nombre qui ayt lesd\textsuperscript{t}. partyes, soyent par exemple vnze partyes données, adioustez 1 pō auoir 12, trouuez des nombres 1\textsuperscript{ers} ou des dignitez qui se multiplians facent 12, ce sont 3 et 4, ostez de chascun l'vnité, reste 2 et 3, exposant d'vn quarré et d'vn cube, donc vn quarré multipliant vn cube, quelz qu'ilz soyent donneront vn nombre qui aura 11 partyes aliquotes, et si l'on prend les moindres nombres pō quarrés et cubes, on aura le moindre nombre de tous ceux, qui ont lesd\textsuperscript{t}. 11 partyes, or le moindre cube est 8, et le moindre quarré est 9 (car il ne faut pas prendre 4, parce qu'il est de mesme rang ou analogie que 8) or 9 par 8 donne 72, qui est le moindre qui ayt lesd\textsuperscript{t}. 11 partyes.
\note{Les chiffres 1 et 2 sont dans la marge gauche. Plus haut dans cette marge, contre le pli, un « 154 » isolé : il fait face à la rangée 21 du carré de 22 de la vue 45 et donne le nombre qu'une tache y couvre (voir la note de cette vue).}

Quand auec vne puissance, l'on trouue vn nombre premier, il ne vaut qu'vn, par exemple soit 20, qui soit fait par 4 et 5, ostez vn de chacun reste 3. et 4, qui est a dire vn nombre impair tel qu'on voudra, et 2 exposant de tel quarré qu'on voudra, donc 3 multipliant 4 (parce qu'on ne peut prendre vn moindre nombre impair, joint au moindre quarré) donnent 12; ou bien 5 multipliant 4 donnent 20, ou 5 multipliant 9 donne 45, et cæt. à l'infiny qui tous ont 5 parties aliquotes seulem\textsuperscript{t}.
\note{Sous « impair », souligné, un petit mot est écrit en dessous, lu « \uncertain{cube} ».}

Et 20 estant donné, on trouue combien il a de partyes aliquotes, car ayant trouué 4 et 5 qui le refont, leurs opposants sont 2 et 1, ausquels l'vnité adioustée l'on a, 3 et 2, qui
\note{La phrase se poursuit en tête de la vue 47.}

\page{47}
\note{Page de gauche (verso) ; aucun chiffre en tête. La première ligne continue la dernière de la vue 46. Les fins de ligne passent dans le pli ; au-delà du pli paraissent des débuts de ligne de la vue 48, qui ne sont pas de cette page.}

se multiplians font 6, dont l'vnité estant ostée reste 5, pō les p\textsuperscript{ties} aliquotes ou bien par la seconde methode, il faut dire 2 fois 1 font 2, et puis 2 et 1 font 3, or trois et deux font cinq.

De mesme pō trouuer les partyes de 360, je prens trois nombres, parce qu'on n'en peut trouuer deux, qui soyent puissances ou premiers, ou premiers et puissances co\textsuperscript{e} il est requis, car l'vnité n'entre point en ces partyes, je prens donc 5, 8, et 9 qui multipliez font 360, or les analogues de 5, 9 et 8 sont 1, 2, et 3 j'adiouste a chascun 1, j'ay 2, 3 et 4, qui multipliez font 24, dont 360 a 24 partyes moins vne, ou bien je dis \struck{et} 1 fois 2 font 2, et 2 et 1 font trois auec deux, ce sont 5 deux fois font 6, 6 et 5 font 11, 11 et 1 font 12 qui auec 11 font 23, mais la 1\textsuperscript{re} façon, qui adiouste 1 à chascun vault mi\add{eux}
\note{« ce sont 5 deux fois font 6 » : lu tel quel ; le calcul de la seconde méthode n'est pas complet sur la page.}

pō trouuer combien 400 a de p\textsuperscript{ties}, je prens 16 et 25, asçauoir vn quarréquarré, et vn quarré dont les exposants sont 4 et 2 qui se multiplians font 8 et par addition 6, c'est a dire en tout 14 autant qu'il a de partyes aliquotes, ou j'adiouste à 2 et 4 l'vnité pō auoir 3 et 5 qui se multiplians font 15 dont ostant l'vnité reste 14 co\textsuperscript{e} deuant

Et pō trouuer vn nombre qui ayt 14 p\textsuperscript{ties}, j'adiouste 1 pō auoir 15, je trouue que 3 par 5 fait 15, j'oste l'vnité de 3 et 5 reste 2 et 4, partant vn nombre quarré et vn quarré quarré donne le nombre, le moindre quarréquarré \uncertain{est} 16, et le moindre quarré est 9.

Item pō trouuer le nombre des p\textsuperscript{ties} de 5000, il faut prendre autant de puissances de 2 et de 5 co\textsuperscript{e} il y a de zero, par exemple icy, 3, et 3, et parce qu'il reste 5 il faut prendre la 4\textsuperscript{me} puissance de 5, de sorte qu'en adioustant 1 à chasque exposant asçauoir à 3 et 4, nous auons 4 et 5 qui multipliez font 20, dont vn estant osté reste 19 parties aliquotes qu'a 5000. S'il y a 15000, il faut prendre trois puissances de 2 et de 5 et à 5 adiouster le 5 co\textsuperscript{e} deuant pō auoir sa 4\textsuperscript{me} puissance reste \ill{} parce que 3 fois 5 font 15, de sorte que nous auons 3, 4, 3, ausquels l'vnité adioustée nō auons, 4, 5 et 4, qui monstrent que 15000 a 80 p\textsuperscript{ties} Aliquotes moins vne.
\note{Après « puissance reste », la fin de la ligne est dans le pli.}

Si vn nombre auoit cent zeros, il faudroit prendre la centiesme puissance de 2 et de 5 et ainsy a l'infiny.

\marginal{3}
L'on sçaura quel nombre a la so\textsuperscript{e} donnée dans ses p\textsuperscript{ties} adioustées ensemble en ceste fa\add{çon} soit la so\textsuperscript{e} 12, je trouue qu'il n'y a que 121, qui ayt ceste proprieté, en ostant vn de douz\add{e} car puis qu'il reste vn nombre premier, asçauoir 11, ou tel autre qui se rencontrera, son quarré donne le nombre cherché.

soit le nombre 24, j'oste 1, reste 23 nombre premier qui multiplié par soy mesme donne vn nombre qui a 24 pō la so\textsuperscript{e} de ses parties aliquotes

soit 25, j'oste vn, reste 24, il faut chercher de quel nombre premier il est composé, je trouue 19, et 5, lesquels multipliez l'vn par l'autre, donnent le nombre 95, dont la so\textsuperscript{e} des parties fait 25.

soit la so\textsuperscript{e} 28, je prends ses costés, qui sont 4 et 7, qui auec leur partyes donnent 7 et 8, qui multipliez auec leurs parties font 56, dont le nombre 28 estant osté reste 28 pō le nombre qui a 28 pō la so\textsuperscript{e} de ses parties aliquotes. Mais il faut en ce cas oster 7 de 28, et non seulement vn, et reste 21. et de certains nombres oster 13 \ill{} ce qui arriue, lors qu'on ne trouue point de nombres premiers qui par addition \ill{} le nombre cherché moins 1.
\note{« so\textsuperscript{e} » : « somme », abrégé comme dans la batch précédente. « 95 » : la somme des parties aliquotes de 95 est 25, comme le dit la note. Les deux \ill{} sont des fins de ligne cachées dans le pli.}

\marginal{4}
vn nombre estant donné, sçauoir la so\textsuperscript{e} de ses p\textsuperscript{ties} aliquotes. soit 12. je prens les nombres premiers ou analogues qui le composent, sçauoir 3 et 4, et dis 3 et ses p\textsuperscript{ties} font 4, 4 et ses p\textsuperscript{ties} font 7, or 4 fois 7 donnent 28, qui contient 12 et la so\textsuperscript{e} de ses p\textsuperscript{ties} ostez 12 reste 16, pō la so\textsuperscript{e} de ses p\textsuperscript{ties} aliquotes.

soit 20, les nombres composans sont 4 et 5 qui auec la so\textsuperscript{e} de leurs parties font 7 et 6, qui
\note{Les mots « 7 et 6, qui » sont écrits seuls au bas de la page, à droite ; la phrase se poursuit à la vue 48. Les chiffres 3 et 4 sont dans la marge gauche.}

\page{48}
\note{Page de droite (recto), chiffrée « 23 » en haut à droite, de la même main fine. La première ligne continue la vue 47 (« 7 et 6, qui »). Au bord gauche de la vue paraissent les fins de ligne de la vue 47.}

7 et 6, qui multipliez font 42, ostez 20, reste 22 pō la so\textsuperscript{e} des p\textsuperscript{ties} de 20

soit 5040, ses p\textsuperscript{ties} composantes, ou nombres premiers, ou leurs puissances sont 16, 9, 7, 5, car multipliez entre eux ilz font 5040, or 16 et ses p\textsuperscript{ties} font 31, 9 et ses p\textsuperscript{ties} font 13, 7 et ses p\textsuperscript{ties} font 8, 5 et ses p\textsuperscript{ties} font 6. Atqui, 31, 13, 8 et 6 se multiplians font 19344, duquel 5040 estant osté reste 14304 pō la so\textsuperscript{e} des p\textsuperscript{ties} aliquotes de 5040.

par mesme moyen l'on trouuera la so\textsuperscript{e} des parties aliquotes de 560400 asçauoir 1374480, et de tel autre nombre qu'on voudra.
\note{« 560400 » : le deuxième chiffre, bouclé, est lu 6 ; un 0 n'est pas exclu. Aucune des deux lectures ne donne 1374480 pour somme des parties aliquotes ; les nombres sont transcrits tels qu'ils sont écrits.}

soit 666.; les nombres dont il est composé, sont 2, 37, 9, eux et leurs p\textsuperscript{ties} font 3, 38, et 13. qui multipliez par eux mesmes font 1282, dont la so\textsuperscript{e} 666 estant ostée, reste 697 pō la so\textsuperscript{e} de ses p\textsuperscript{ties} qui sont 11 en nombre, car les exposants de 2, 9, et 37 sont 1, 2 et 1 ausquels l'vnité adioustée, l'on a 2, 3, 2 lesquels multipliez entre eux, l'on a 12, dont il faut oster 1 reste 11
\note{« 1282 » et « 697 » sont lus tels quels ; le produit $3 \cdot 38 \cdot 13$ vaut 1482, et $1482 - 666 = 816$ ; $1282 - 666$ vaut 616. Le calcul reste tel qu'il est écrit.}

soit 12000. ses p\textsuperscript{ties} composantes sont 3, 125, 32. Elles auec leurs p\textsuperscript{ties} font, 4, 156, 63 qui multipliez entre elles font 39312, dont la so\textsuperscript{e} estant ostée asçauoir 12000, reste, 273\uncertain{1}2 pō la so\textsuperscript{e} de ses p\textsuperscript{ties}. Et pō sçauoir combien il a de p\textsuperscript{ties}. je prens, les exposans de 3, 32 et 125, qui le composent asçauoir 1, 5 et 3 ausquelz l'vnité estant adioustée, j'ay 2, 6, et 4, qui multipliez par eux mesmes, font 48, dont l'vnité estant ostée, il reste 47 pō le nombre des p\textsuperscript{ties} de 12000.
\note{« 273\uncertain{1}2 » : le quatrième chiffre est repassé, 1 sur 2 ou l'inverse ; $39312 - 12000 = 27312$.}

soit 5000 les 3 zero donnent la 3\textsuperscript{me} puissance de 2 qui est 8, et la 3\textsuperscript{me} puissance de 5, à laquelle 5 adiousté, nō auons la 4\textsuperscript{me} puissance de 5, asçauoir 625, qui multiplié par 8 refait le nombre 5000, les exposants desd\textsuperscript{t}. puissances, sont 3 et 4, l'vnité leur estant adioustée, nō auons 4 et 5, qui multipliez font 20, ostez en vn, reste 19 pō le nombre de ses p\textsuperscript{ties}. la so\textsuperscript{e} de ses p\textsuperscript{ties} se trouue en prenant 8 et 625 auec leur p\textsuperscript{ties} asçauoir 15 et 781 qui multipliez, donnent 11715, duquel la so\textsuperscript{e} ostée de 5000 reste 6715 pō la so\textsuperscript{e} de ses p\textsuperscript{ties}

soit le nombre 888, qui respond a 1468\uncertain{5}. Il est composé de 8, 3, et 37 dont les exposants sont 3, 1, et 1 ausquelz l'vnité adioustée nous auons 4, 2, et 2, qui multipliez donnent 16, dont 1 osté reste 15 pō le nombre de ses p\textsuperscript{ties}, dont la so\textsuperscript{e} se trouue en prenant 8 et ses parties qui font 15, et 3 et ses parties, qui font 4 et 37 et ses p\textsuperscript{ties} qui font 38 qui multipliez font 2280, dont il faut oster 888, et l'on a 1392 pō la so\textsuperscript{e} de ses p\textsuperscript{ties}
\note{« qui respond a 1468\uncertain{5} » : le dernier signe peut être un 5 ou un s ; le sens de la phrase n'est pas établi.}

le moindre nombre de ceux qui contiennent 24 fois et non plus trois nombres proportionaux, est 24843 et non 753571, qui les contient aussy. pō auoir tant de fois 3 nombres proportionaux qu'on voudra, il faut doubler le nombre cherché, par exemple, je veux le moindre qui en contienne 4, je double 4, j'ay 8, j'adiouste 1 fait 9, qui est composé de 3 et 3, ostez 1, reste 2 et 2, et leur moitié est 1 et 1 qui representent ou exposent 2 nombres premiers, qui doibvent estre 6 plus vn, ou multiples de 6+1 soyent donc 7 et 13, qui multipliez font 91, dont les quatre trinitez de proportionaux 1.9.81 :| 25.30.36 :| 13.26.52 :| 7.21.63 :| qui tous refont 91.
\note{Au bout de la dernière ligne, dans la marge droite, un signe en forme de double chevron, peut-être un renvoi.}

\marginal{Si apres auoir doublé le nombre qu'on veut de nombres proportionaux \ill{} 3 à 3, et y auoir adiousté 1, il se trouue \add{im}pair, co\textsuperscript{e} quand on cherche vn nombre qui en contienne 2 seulem\textsuperscript{t}, l'on double 2 auec l'vnité adioustée l'on a 5, il faut seulem\textsuperscript{t} prendre le quarré de 7, qui est 49, parce que deux que l'on cherchoit sont l'exposant: les 2 ternaires de nombres proport., sont 9.15.25: 7.14.28/ si l'on multiplie 49 par 3, le produit contiendra \ill{} ternaires proportio. mais tel autre nombre qu'on voudra co\textsuperscript{e}, 2, 4, \ill{}, et cæt. multipliant n'augmente point lesd\textsuperscript{t}. trinitez de nombres proportionels et mesme le ternaire ne l'augmente que la seule premiere fois, car s'il multiplie encore vne fois, 2 fois et cæt. n'y adiouste plus rien/}
\note{Cette note est écrite dans la marge gauche, contre le pli, en face des deux derniers paragraphes ; ses débuts de ligne sont cachés dans le pli, et les \ill{} en sont. Ses deux dernières lignes s'étendent sous le texte, sur toute la largeur de la page. Un mot de la marge, « prendre », est souligné d'un trait qui rejoint le texte.}

vn nombre estant donné, on cognoist s'il n'est pas premier

1\textsuperscript{o} s'il finist par 5, ou par 0, car pō lors il est diuisible par 5.

2\textsuperscript{o} si en ostant tous les 9, il reste 9, 6, ou 3 il n'est pas premier par exemple 111 est diuisible par 3 parce qu'il n'y a que 3 vnités. 75512713527351 n'est pas premier, car il ne reste rien apres 9, donc il est diuisible par 9.

Soit 6000 nombre des années de la durée du monde. ses p\textsuperscript{ties} aliquotes sont 39 en nombre, qui se trouuent par les exposantz de 3, 16, 125. asçauoir, 1, 4, 3 car y adioustant vn, on a 2, 5, 4 qui font 40, osté l'vnité, reste 39 pō les p\textsuperscript{ties} aliquotes. lesd\textsuperscript{t}. nombres 3, 16, 125, qui sont les analogues refaisant par multiplicaõn 6000, auec leurs parties aliquotes à sçauoir, 4, 31, et 156 qui se multiplians, font vne somme qui comprend celle des p\textsuperscript{ties} aliquotes de 6000, et de plus 6000, lequel osté reste, 13344 pō la so\textsuperscript{e} des partyes aliquotes de 6000 duquel vous verrez lesd\textsuperscript{t}. 39 parties aliquotes a la page suyuante/

\page{49}
\note{Cette vue montre un petit feuillet volant, de papier plus clair, coupé en biais à son angle inférieur droit, posé sur la page de gauche de l'ouverture suivante et en cachant la moitié ; la page elle-même est lue à la vue 50, où le feuillet a été écarté. Rien n'indique à quelle page du livre ce feuillet se rapporte. Il est de la même main que le reste. Les combinaisons sont écrites en colonnes, chacune entre accolades ; elles sont transcrites ici ligne par ligne, dans l'ordre du feuillet : la somme des deux dés, les coups qui la donnent (dans les deux ordres) et leur nombre, puis les coups sans l'ordre et leur nombre, puis les petits traits de décompte de la troisième colonne, notés « 1– » ou « 1+ » selon leur forme.}

pō deux dez

2 : 1 1, 1 1 — 2 \quad 1 1 — 1 \quad 1+

3 : 2 1, 1 2 — 2 \quad 2 1 — 1 \quad 1–

4 : 2 2, 2 2, 3 1, 1 3 — 4 \quad 3 1, 2 2 — 2 \quad 1–, 1+

5 : 4 1, 1 4, 3 2, 2 3 — 4 \quad 4 1, 3 2 — 2 \quad 1–, 1–

6 : 5 1, 1 5, 4 2, 2 4, 3 3, 3 3 — 6 \quad 5 1, 4 2, 3 3 — 3 \quad 1–, 1–, 1

7 : 6 1, 1 6, 5 2, 2 5, 4 3, 3 4 — 6 \quad 6 1, 5 2, 4 3 — 3 \quad 1–, 1–, 1–

8 : \struck{7 1}, \struck{1 7}, 6 2, 2 6, 5 3, 3 5, 4 4, 4 4 — \struck{8} 6 \quad \struck{7 1}, 6 2, 5 3, 4 4 — \struck{4} 3 \quad 1–, 1–, 1+

9 : \struck{8 1}, \struck{1 8}, \struck{7 2}, \struck{2 7}, 6 3, 3 6, 5 4, 4 5 — \struck{8} 4 \quad \struck{8 1}, \struck{7 2}, 6 3, 5 4 — \struck{4} 2 \quad 1–, 1–

10 : \struck{9 1}, \struck{1 9}, \struck{8 2}, \struck{2 8}, \struck{7 3}, \struck{3 7}, 6 4, 4 6, 5 5, 5 5 — \struck{10} 4 \quad \struck{9 1}, \struck{8 2}, \struck{7 3}, 6 4, 5 5 — \uncertain{\struck{5} 2} \quad 1–, 1–

11 : \struck{10 1}, \struck{1 10}, \struck{9 2}, \struck{2 9}, \struck{8 3}, \struck{3 8}, \struck{7 4}, \struck{4 7}, 6 5, 5 6 — \struck{10} 2 \quad \struck{10 1}, \struck{9 2}, \struck{8 3}, \struck{7 4}, 6 5 — \uncertain{\struck{5} 1} \quad 1–

12 : 6 6, 6 6 — 2 \quad 6 6 — 1
\note{Le copiste a d'abord écrit, pour les sommes de 8 à 11, toutes les décompositions en deux nombres, puis il a barré d'un trait horizontal celles où entre un nombre supérieur à 6, et corrigé les totaux (8 en 6, 8 en 4, 10 en 4, 10 en 2 ; 4 en 3, 4 en 2 dans la seconde colonne). Pour 10 et 11, le total de la seconde colonne est un 5 suivi d'un second chiffre en exposant, à demi caché sous le bord de la page de la vue 49 ; il est lu « 2 » et « 1 », mais la biffure du 5 n'est pas sûre. Pour 4 et 6, le deuxième « 2 2 » et le deuxième « 3 3 » ont un de leurs chiffres repassé. Pour 2, le copiste compte « 1 1 » deux fois, comme pour 4, 6, 8, 10 et 12 les coups doubles ; le « 1 » du premier est une sorte de T. Un « 1 » isolé, peut-être biffé, précède la seconde accolade de la somme 6.}

\page{50}
\note{Page de gauche (verso) ; aucun chiffre en tête. C'est la page sur laquelle était posé le feuillet volant de la vue 49, ici écarté. Au bord droit, au-delà du pli, paraissent une colonne de chiffres (2 à 12, avec « … 18 », « … 17 »…) et le bord d'un autre feuillet, qui appartiennent à la vue 51.}

\[
\begin{array}{|r|r|r|}
\hline
1 & 6 & 8 \\
5 & 30 & 40 \\
25 & 150 & 200 \\
125 & 750 & 1000 \\
3 & 4 & 24 \\
15 & 20 & 120 \\
75 & 100 & 600 \\
375 & 500 & 3000 \\
2 & 12 & 16 \\
10 & 60 & 80 \\
5 & 300 & 400 \\
120 & 1500 & 2000 \\
 & & 48 \\
 & & 240 \\
 & & 1200 \\
\hline
\end{array}
\]
\note{Ce tableau est dans l'angle supérieur gauche, sous le titre « 39 partyes aliquotes de 6000 », précédé d'un petit trait. Il donne trente-neuf nombres : les parties aliquotes de 6000 qu'annonçait la fin de la vue 48 (« a la page suyvante »). Deux nombres de la première colonne sont lus « 5 » et « 120 » ; la suite des autres (1, 5, 25, 125 ; 3, 15, 75, 375 ; 2, 10…) appellerait 50 et 250, et 5 et 120 sont déjà ailleurs dans la table ; ils sont transcrits tels qu'ils sont écrits.}

6679 n'est pas premier, il se diuise par 23 et 29. mais 1373 est premier et non 1369 quarré de 37, ny \struck{\ill{}} 1371

viã coĩ nõ reperiz, an 1073602561, sit primus necne nisi per \struck{\ill{}} oẽs primos vsque ad 8191 diuiseris.
\note{« viã coĩ nõ reperiz » : lu tel quel ; par inférence « viam communem non reperies ». « oẽs » : « omnes ».}

Numeri triangularis imparis summa habetur sumpto medio numero et ducto in \struck{\ill{}} \add{praedictum} numerum, vt triangularis 7, habet pro medio 4, inter 3 et 3, qui 4 ductus in 7 dat 28 pro summa partium trianguli, cuius latus vel basis 7. si par fuerit vt 6, media pars 3 ducta in seq\textsuperscript{ti} trianguli 7 dat 21 pro summa trianguli 6.
\note{« praedictum » est écrit au-dessus du mot biffé ; il est mis ici en \add{} parce que l'abréviation (« pradictu ») est développée. « seq\textsuperscript{ti} » : lu comme l'abréviation de « sequenti(s) » ; la lecture « sexti » n'est pas exclue.}

Numerus omnium praecedentium triangulorum habetur, si sumatur \struck{a} $\frac{1}{3}$ seq\textsuperscript{tis} trianguli et illud $\frac{1}{3}$ ducatur in praecedentis trianguli latus. Sint 5 primi trianguli ab vnitate inclusiue, sumatur $\frac{1}{3}$ \uncertain{sexti} trianguli hoc est 7 cum \uncertain{sextus} triangulus habeat 21 vnitates at 7 ductus in 5 dat 35 pro summa triangulorum ab vnitate inclusiue.
\note{« in praecedentis trianguli latus » : « latus » est au bout de la ligne, contre le pli, et n'est lu que par le sens. « sexti », « sextus » : un s long suivi d'une lettre bouclée, qui peut être le q de « seq. » ; le sixième triangle vaut bien 21.}

Si quaeratur summa 12 triangulorum inclusiue à 13\textsuperscript{o} triangulo, sumenda est summa 24 triangulorum qua habita, et dempta summa 12 primorum, supererit summa quaesita.

Numeri ciuilis platonici 5040 partes aliquotae \struck{54} \textsuperscript{59}. Illius p\textsuperscript{tes} componentes 9, 16, 7, 5 sequuntur p\textsuperscript{tes} nec vllus minor eas habet
\note{« 54 » est biffé d'un trait épais et « 59 » écrit au-dessus.}

\[
\begin{array}{|l|r|r|r|r|r|r|r|r|r|r|r|}
\hline
1 & 3 & 9 & 5 & 15 & 45 & 7 & 21 & 63 & 35 & 105 & 315 \\
2 & 6 & 18 & 10 & 30 & 90 & 14 & 42 & 126 & 70 & 210 & 630 \\
4 & 12 & 36 & 20 & 60 & 180 & 28 & 84 & 252 & 140 & 420 & 1260 \\
8 & 24 & 72 & 40 & 120 & 360 & 56 & 168 & 504 & 280 & 840 & 2520 \\
16 & 48 & 144 & 80 & 240 & 720 & 112 & 336 & 1008 & 560 & 1680 & \\
\hline
\end{array}
\qquad
\begin{array}{|c|r|}
\hline
1 & 5040 \\
\hline
2 & 2520 \\
2 & 1260 \\
2 & 630 \\
2 & 315 \\
3 & 105 \\
3 & 35 \\
5 & 7 \\
\hline
\end{array}
\]
\note{Entre les deux tableaux, les mots « methodus reperiendi partes ». Dans le premier, les nombres de chaque colonne sont écrits un peu décalés vers le bas d'une colonne à l'autre (les colonnes 5, 15, 45… commencent une ligne plus bas dans l'écriture) ; ils sont alignés ici. Le second tableau est la décomposition de 5040 par divisions successives.}

\marginal{1}
Lapides distant à \ill{} per vnum passum singuli, v.g. sint 20 lapides, scitur numerus omnium passuum duobus modis, vel ducto 20 in 21, vel quadrato 20 addatur numerus lapidum.
\note{Un blanc est laissé dans la ligne après « à », comme au paragraphe suivant et au paragraphe 4 ; le copiste semble n'avoir pas lu un mot de son modèle.}

\marginal{2}
Si lapis primus distet à \ill{} 1 passu, 2\textsuperscript{us} lapis à priori duobus passibus, tertius a 2\textsuperscript{o} passibus tribus et sic deinceps, sint lapides centum, scientur passus, si ducatur $\frac{1}{3}$ trianguli seq\textsuperscript{tis} (numero 1717 in 100, hinc fit 171700 pro dimidia p\textsuperscript{te} passuum, ergo erunt passus 343400. si numerus lapidum per 3 diuidi possit, per illius numeri $\frac{1}{3}$ triangularis numerus sequens multiplicandus.
\note{« (numero » : la parenthèse ouverte n'est pas fermée ; le mot qui suit est lu « numero » avec doute, un mot abrégé en « nump° ». $\frac{1}{3}$ du 101\textsuperscript{e} triangle ($5151$) vaut 1717.}

\marginal{3.}
Si lapides 12 sequantur analogiam seu progressionem 2, vt primus distet 1, 2\textsuperscript{us} 4, 3\textsuperscript{us} 8 et cæt passus faciendi reperientur, si 12\textsuperscript{a} potestas \struck{\ill{}} numeri 2 (4096) sumatur, cuius duplum 8192 a quo dempto \textsuperscript{14} numero lapidum + 2 supererit 8178 quae est media pars passuum, ergo erunt passus 16356./
\note{Le « 14 » est écrit au-dessus de « numero lapidum + 2 », soit 12 + 2.}

\marginal{4.}
Si sumantur 5 lapides, et is à \ill{} distet 1, 2\textsuperscript{us} à 1\textsuperscript{o} 3, 3\textsuperscript{us} à 2\textsuperscript{o} 9 et cæt. juxta progressionem geometricam 3, reperietur summa passuum si sumatur 5\textsuperscript{a} potestas 3 nempe 243 (nam lapidum numerus est exponens potestatum) aufer 1 ab 243 et 3 qui est radix analogiae supersunt 242 et 2 diuido 242 per 2, quotiens 121, ostendit maximum spatium a \uncertain{vrbe} ad vltimum lapidem. habetur autem summa omnium spatiorum ducto 3 in 121. fit 363, ex quo numerus lapidum ablatus, superest 358, qui diuisus per 2 hoc est radice analogiae +1, exurgit 179 pro media p\textsuperscript{te} passuum, sic erunt passus. 358.
\note{« et is » : lu tel quel (« et 1\textsuperscript{us} » ?). « a urbe » : lu « a urbe » ou « a verbo », sous l'abréviation « a vrbe ». Les nombres sont transcrits tels qu'ils sont écrits : $242/2 = 121$, $3 \cdot 121 = 363$, $363 - 5 = 358$, $358/2 = 179$, et le copiste conclut « passus. 358 ».}

\marginal{5}
Si non jubet se sedere à corbe sumant suas differentias, siue praecedentis lapidis a qu\ill{} tempore afferantur in corbem, ac sint in 4\textsuperscript{o} praecedenti puncto lapis vltimus. Igitur si hic fuerint 5 lapides, fient 242 passus pro oĩb\textsuperscript{9} lapidibus afferend\add{is} qui est duplus 121, quibus distabat praecedens vltimus lapis a corbe.
\note{Paragraphe peu sûr dans sa lettre : « non jubet se sedere » et « sint in 4\textsuperscript{o} praecedenti puncto » sont lus tels qu'ils sont écrits, sans que le sens soit établi ; « corbe », « corbem » sont nets. « oĩb\textsuperscript{9} » : « omnibus ».}

Cuborum omnium praecedentium numerus scitur, sumendo quadratum trianguli maioris cubi, vt sint 30 cubi, triangulus 30 sit 465, cuius quadratum 216225 dat summam 30 primorum cuborum ab 1 incipientium. Contingit vero vt in 4 primis 1.8.27.64 triangularis numerus quadrati 4, hoc est 16 (qui est 136) dat summam \uncertain{summarum} \ill{} summam 100 pro 4 primis cubis junctam summae 3 priorum 36

Summa quadratorum habetur, vt 4 primorum sumpto trigono sequentis 5. id est 15. cuius $\frac{1}{3}$ ductum in 4 facit 20, qui duplicatus facit 40, a quo demptus triangulus 4

\note{Réclame au bas de la page : « 10, supersu… », au bord du pli ; la phrase se poursuit en tête de la vue 51.}
\note{Dans la marge droite, contre le pli, une colonne de nombres de 1 à 12, chacun suivi de points : ils sont de la vue 51 (voir cette vue).}

\page{51}
\note{Page de droite (recto), chiffrée « 24 » en haut à droite, de la même main fine. La première ligne continue la réclame de la vue 50 (« 10, supersu… »). Les deux tables de dés sont écrites dans la marge gauche, contre le pli, en face des paragraphes « Des differens ieux des dez » ; elles sont données à la fin de la vue. Au bord gauche de la vue paraissent les fins de ligne de la vue 50.}

10 supersunt 30 pro summa 4\textsuperscript{or} primorum quadratorum. 1.4.9.16: sint 30 prima quadrata, sume trianguli 31 id est 496; qui quia non habet $\frac{1}{3}$ absque fractione vt vites fractionem, sume $\frac{1}{3}$ 30 id est 10 qui ductus in 496 dat 4960, qui duplicatus dat 9920, à quo demptus triangulus 30 \uncertain{id est} 465 superest 9455, pro summa 30 primorum quadratorum
\note{« sume trianguli 31 » : lu tel quel, sans « triangulum ». Le copiste abrège « id est » par « id ẽ » ; devant 465, le signe est lu de même, avec doute.}

Quarrez quarrez ou \struck{4\textsuperscript{mes} p\textsuperscript{ce}} 4\textsuperscript{mes} puissances et leurs differences

La difference des cubes se suyuants immediatem\textsuperscript{t} de 1, 2, 3, 4, et cæt. different par les hexagones centraux, co\textsuperscript{e} de 1 à 8, il y a 7, et 6 est le p\textsuperscript{r} hexagone central adioustant 1, c'est la difference, de mesme de 8 à 27, il y a 18 plus vn, or 18 est le triangle de \struck{\ill{}} 2 c'est à dire 3, multiplié par 6: et 6 sert à multiplier tous les aut\textsuperscript{res} triangles, co\textsuperscript{e} entre 27 et 64, multipliez le triangle de 3 c'est a dire 6, par 6, vous aurez 36 qui, plus vn, est la difference de 27 a 64.
\note{« et 6 est le p\textsuperscript{r} hexagone central adioustant 1 » : lu tel quel.}

Mais la difference des quarrez quarrez se trouue adioustant la colomne massiue de l'hexagone central collateral au plus grand quarré quarré à la colomne quarrée massiue (qui est le cube) collateral au moindre quarré quarré. Ainsi pō auoir la difference du quarréquarré de 3 et de celuy de 4 prenez la colonne hexagone de 4 qui est 148, et l'adioustez au cube de 3 qui est 27, pō auoir 175
\note{« à » est souligné sur la page. « 148 » : lu tel quel.}

Autrem\textsuperscript{t}. adioustez ensemble les colonnes creuses des hexagones centraux, ou bien multipliez la pyramide quarrée vuide collaterale au plus grand quarré quarré par l'impair qui est en mesme qualité ou bien multipliez l'hexagone central collateral au plus grand quarré quarré par son costé; et adioustez au produit la so\textsuperscript{e} des hexagones precedens.

Q. C. ou 5\textsuperscript{mes} puissances et leurs differences.

prenez le double du triangle collateral à la moindre puissance, et puis le triangle de ce double, et mettez 1 a la fin du triangle co\textsuperscript{e} pō la difference des 5\textsuperscript{mes} puissances de 3 et de 4, prenez le triangle de 3 qui est 6, son double 12 dont le triangle est 78, et puis 1 vous aurez 781.

Des differens ieux des dez.

Tous les ieux de trois dez, qui peuuent venir sont le cube de six, c'est à dire 216. mais les ieux particuliers, c'est à dire les diuerses façons de chacun des nombres, on la en cette façon, le moindre est 3, et le plus grand 18, qui sont en tout 16 nombres, dont les 8 premiers ressemblent aux 8 derniers, et ainsy 3 et 18 se raportent 4 et 17 | 5 et 16 | \textsuperscript{6 et 15.} 7 et 14 | 8 et 13, 9 et 12 / 10 et 11. pō les diuerses façons dont on peut auoir chascun desd\textsuperscript{t}. nombres, il faut pō les six premiers ou les six derniers, prendre les 6 premiers triangles, et ainsy on pourra auoir 3 et 18 en vne sorte, 4 et 17 en trois sortes, 5 et 16 en six sortes, 6 et 15 en 10 sortes, 7 et 14 en 15 sortes, 8 et 13 en 21 sortes. pō les 2 suyuants 9 et 12, il faut prendre le quarré de 5, sçauoir 25. et pō 10 et 11 on prendra le cube de trois, sçauoir 27 or si l'on assemble ces nombres 1, 3, 6, 10, 15, 21, 25, 27 on aura 108, dont le double est 216.
\note{« on la en cette façon » : lu tel quel. « 6 et 15. » est écrit au-dessus de la ligne, entre « 5 et 16 » et « 7 et 14 ». Après les six premiers triangles, 25 et 27 rompent la suite ; la somme 108 et son double 216 sont justes.}

pō les ieux de deux dez ilz peuuent arriuer de 62 façons et pourtant il n'y a que 31 \struck{\ill{}} ieux differents/ car 2 et 3 poincts peuuent arriuer chascun de deux façons, quatre et cinq de 4 façons, six et sept de 6: 8 et 9 de 8 façons, 10 et 11 de 10 et 12 de deux. mais il ne faut prendre que la moitié de chascun;
\note{Les nombres (62, 31, « 8 et 9 de 8 façons », « 10 et 11 de 10 ») sont transcrits tels qu'ils sont écrits ; la table de la marge, plus bas, compte autrement (42 et 21).}

Les 1\textsuperscript{ers} quarrez

1 : 16 . 81 : 256 625 2401 \quad leurs differences 15 . 65 . 175 . 369 . 1776
\note{Sous chaque intervalle de la suite 1, 16, 81, 256, 625, 2401, la différence est écrite sous un petit arc : 15, 65, 175, 369, 1776. La suite saute 1296 (six à la quatrième) : 1776 est la différence de 625 à 2401.}

\marginal{pō trois dez}
\[
\begin{array}{rcr||rr}
3 & \cdots & 18 & \cdots\ 1 & \cdots\ 1 \\
4 & \cdots & 17 & \cdots\ 3 & \cdots\ 1 \\
5 & \cdots & 16 & \cdots\ 6 & \cdots\ 2 \\
6 & \cdots & 15 & \cdots\ 10 & \cdots\ 3 \\
7 & \cdots & 14 & \cdots\ 15 & \cdots\ 4 \\
8 & \cdots & 13 & \cdots\ 21 & \cdots\ 5 \\
9 & \cdots & 12 & \cdots\ 25 & \cdots\ 6 \\
10 & \cdots & 11 & \cdots\ 27 & \cdots\ 6 \\
\hline
 & & & 108 & 28 \\
 & & & 2 & \\
\hline
 & & & 216 &
\end{array}
\]

\marginal{pō 2 dez.}
\[
\begin{array}{rcrr}
2 & \cdots\ 12 & 2 & \cdots\ 1 \\
3 & \cdots\ 13 & 2 & \cdots\ 1 \\
4 & \cdots\ 14 & 4 & \cdots\ 2 \\
5 & \cdots\ 15 & 4 & \cdots\ 2 \\
6 & \cdots\ 16 & 6 & \cdots\ 3 \\
7 & \cdots\ 17 & 6 & \cdots\ 3 \\
8 & \cdots & 6 & \cdots\ 3 \\
9 & \cdots & 4 & \cdots\ 2 \\
10 & \cdots & 4 & \cdots\ 2 \\
11 & \cdots & 2 & \cdots\ 1 \\
12 & \cdots & 2 & \cdots\ 1 \\
\hline
 & & 42 & 21
\end{array}
\]
\note{Les deux tables sont dans la marge gauche, séparées par un trait. Dans celle des trois dés, les chiffres 3 à 10 de la première colonne sont contre le pli, à demi cachés ; les sommes qu'ils appellent (3 + 18 = 21) les confirment. Un signe en forme de crochet, sous la colonne, renvoie au texte. Dans celle des deux dés, la colonne 12 à 17, entre les points, est barrée d'un lavis vertical et n'est pas continuée ; le 1 de « 11 » et le 3 de « 13 » sont peu nets. Les chiffres des lignes 8 à 11 sont repassés à l'encre épaisse : 6 sur 8 à la ligne 8, 4 sur 8 à la ligne 9, 2 sur 12 (?) à la ligne 11, et dans la dernière colonne 3 sur 2 (ligne 8) et 2 sur 3 (ligne 9), avec un petit 3 et un petit 2 récrits en marge. Des accolades réunissent les lignes deux à deux (2–3, 4–5, 6–7, 8–9, 10–11). Le total 42 est écrit sur un autre chiffre, peut-être 4.}

\page{52}
\note{Page de gauche (verso du feuillet chiffré 24), sans texte propre : ce qu'on y voit est la vue 51 par transparence, à l'envers. Seule une ligne pâle au pied de la page se lit à l'endroit : « Ducta \ill{} 4 facit 20 qui duplicatus facit 40 a quo demptus triangulus \ill{} », qui répète presque mot pour mot la dernière ligne de la vue 50. Elle est si effacée qu'on ne peut dire si c'est une ligne commencée puis lavée, ou un report d'encre ; elle n'est pas donnée comme texte.}

\page{53}
\note{Page de droite (recto), chiffrée « 25 » en haut à droite, de la même main fine. Même copiste, d'une écriture plus posée ; le texte est en latin et forme une pièce à part, qui court sur les vues 53 à 55.}

Academia Parisiensis viros Clarissimos Galilaei familiares et amicos lyncaeos precatur, vti sequentibus in dialogorum libros notis respondeant

In primis quae primo dialogo pag. 2\textsuperscript{a} et 3\textsuperscript{a} habet.

1\textsuperscript{o}. Dicimus machinam maiorem confici debere quae sit eo durior, et ruptu difficilior, quo illius pondus et figura maiora fuerint: tantumque fore discrimen inter minorem et maiorem eiusd\textsuperscript{m}. materiae, quantum fuerit inter duas aequalis magnitudinis, quarum vna constat ex materia minus ponderosa, et tamen duriori quam altera.

2\textsuperscript{o} Dicimus non esse vacuum quod impedit ne duae laminae politae sibi superpositae separari queant, sed aerem desuper prementem et se toto incumbentem pag. 23. neque vacuum impedire trombas, seu Ctesibica organa ne trahant aquam pag. 17.

3\textsuperscript{o} Dicimus in discursu figurae pag. 28\textsuperscript{a}. videri esse sophisma seu paralogismum: neque aliud hinc concludi posse nisi lineam aut superficiem, non esse maius solidum quam punctum; minime vero lineam aut superficiem absolute sumptam, non esse maiorem puncto

4\textsuperscript{o} Dicimus specula illa Archimedea pag. 42 esse impossibilia, et in eo Cavalerium ignoratione \uncertain{vitrea} catoptricae errasse, cuius rei rationem explicabimus, si quis eadem re laboret.
\note{« vitrea » : écrit « vira » avec un trait suscrit.}

5\textsuperscript{o} \uncertain{Negamus} pag. 50 circulum spatiola, aut puncta vacua relinquere, dicimusque in eo nullum aliud esse mysterium, nisi quod minor circulus maiore lentius moueatur.
\note{« Negamus » : la capitale ressemble à un L ; le sens demande N.}

6\textsuperscript{o} Quod pag. 54 dicitur de tracto auri filo, nihil facit ad rem, non enim patitur rarefactionem, sed mutat duntaxat figuram.

7\textsuperscript{o} Certum est aquam resistere diuisioni, alioquin non adhaereret digitis tangentibus contra pag. 70.

8\textsuperscript{o} Supponit corpora grauitare in vacuo pag. 73 quod falsum est

9\textsuperscript{o}. Soni chordarum aurearum non sunt grauiores ob auri maius pondus, sed ob maiorem mollitiem. Estque falsum pondus alicuius corporis magis velocitati motus, quam illius crassitiem et magnitudinem resistere contra pag. 103.

10\textsuperscript{o}. Non probat transuersum prisma esse instar vectis cuius hypomochlion sit in medio densitatis, neque id verum est.

11\textsuperscript{o}. Dialogus tertius factus videtur, vt demonstretur omnes vibrationes eiusd\textsuperscript{m} appensae, duratione aequales esse quod tamen minime probat.
\note{Les renvois de cette page (« pag. 2\textsuperscript{a} et 3\textsuperscript{a} », 23, 17, 28, 42, 50, 54, 70, 73, 103) sont aux pages des dialogues de Galilée que la pièce discute, non à l'Harmonie universelle.}

\page{54}
\note{Page de gauche (verso) ; en tête, à gauche, une petite marque « 1 » ou « i », pas un numéro de page. Suite de la pièce latine de la vue 53. Les fins de ligne passent dans le pli ; les lettres cachées ne sont complétées par \add{} que lorsque le mot est certain.}

Et pag. 166. nisi sit verum quod ait de planis inclinatis, tot\add{us} tractatus falsus erit, itaque ea demonstrare hoc opus hic labor es\add{t}

12\textsuperscript{o} Si quae ait pag. 236. à sectione illa, mobile quoddam super planum Et cæt.) vera non sunt, corruit totus tractatus. Atqui vera esse non probat.
\note{La parenthèse fermante n'a pas d'ouvrante.}

13\textsuperscript{o}. Conuersam suae propositionis facit pag. 288 absque demonstrati\add{one} et explicatione, nempe si ictus horizontalis à B, versus E, sequ\ill{} parabolam BD, ictum obliquum DE, eamdem parabolam DB, descripturum.
\note{Dans « parabolam DB », le D est repassé sur une autre lettre.}

14\textsuperscript{o}. Aduersus ea quae pag. 280 dicit de amplitudinibus semi\ill{} parabolarum et cæt. Dicimus haec omnia repugnare experientiae 1\textsuperscript{o} quia hinc sequeretur ictum seu parabolam 45\textsuperscript{e} graduum, esse 28\textsuperscript{es}. ictu maiorem, seu parabola vnius gradus, atqui consta\add{t} experientia illius parabolae, amplitudinem huius ad summum esse triplam. 2\textsuperscript{o} sequeretur bombardam minorem (quam nos arquebusiam appellamus) horizonti parallelam, numquam scopum attingere, sed 12 pedes integros pilam sub scopum descende\add{re} Cum enim spatio vnius secundi minuti pila moueatur de puncto (vt vocant) ad punctum. Sit bombarda AC, quae collineet ad B, sint vero ducenti passus ab: A, ad B. Certum est experientia pila\add{m} à C ad B, moueri spatio secundi minuti vnius, certumque scopum propinquè peteret in B, Et tamen spatio vnius secundi minuti cert\add{um} est pilam versus centrum descendere à B, ad D per spatium duodecim pedum. Itaque si motus horizontalis AB, non impedit motum ad centrum BD, seu AE, quod dixi contingere, contra perpetuam experientiam
\note{Figure dans la marge gauche : un rectangle A B D E, A C B en haut (le canon AC le long du côté supérieur), la diagonale de A à D, les côtés AE et ED en pointillé, BD en trait plein. « 28\textsuperscript{es}. » : lu tel quel (« vicies octies »). « peteret » : lu tel quel.}

15\textsuperscript{o} Quod ait grauia descendentia eo tardius in aquâ, aere Et cæt. descendere, quam in vacuo, quo plus de pondere illorum tollitur ab aere, aquâ, vel alio medio, repugnat experientiae. Quod ita demonstramus. Sit aer 400\textsuperscript{es}. aquâ leuior, plumbum duodecuplo grauius aqua, igitur plumbum 4800\textsuperscript{es} grauius erit aere atque adeo plumbum aeque velociter in aere descendet ac in vacu\add{o} dempta $\frac{1}{4800}$ parte, in aquâ vero tardius quam in vacuo descend\add{et} $\frac{1}{12}$ parte itineris. Igitur plumbum 12 tantum pedes in aere descend\ill{} quandiu in aqua 11. Atqui constat experientia eodem tempor\add{e} quo 12 pedes descendit in aqua, pedes 48 descendere in aere.

Porro rogamus viros clarissimos vt velint nos docere quant\ill{} fuerit pondus, quataeue magnitudo lagenarum quibus Galilae\add{us}

\note{Réclame au bas de la page, à droite : « aeris pondu… », reprise en tête de la vue 55.}

\page{55}
\note{Page de droite (recto), chiffrée « 26 » en haut à droite, le chiffre souligné, de la même main fine que les précédents. Fin de la pièce latine des vues 53 et 54 : la première ligne continue la réclame de la vue 54 (« aeris pondu… »). Les lettres qui paraissent au bord gauche de la vue sont les fins de ligne de la vue 54. La vue 56, verso de ce feuillet, n'a pas de texte : on y voit cette page par transparence et le cachet rond « BIBLIOTHEQUE IMPERIALE MSS. » ; le même cachet paraît, pâle, au milieu du diagramme de cette page.}

aeris pondus explorauit, cuius etiam ponderis et magnitudinis aer ponderi probatus: et quodnam pondus bilances ex aequilibrio deiecit, vt illius obseruationes proxime imitemur. Denique cum tantopere laborarit, in inueniendâ vi percussionis illos vehementer Rogamus vt tractatum illum studiosis non denegent.
\note{« aer ponderi probatus » : lu tel quel.}

Haec sunt quae ipsi Galilaeo etiamnum viuenti atque florenti Parisienses scripserant quorum missionem tanti viri obitus praeuenit, et nobis reliquit luctum illum, quem omnes eruditi testantur; qui cum crediderint non deesse plures illius amicos, qui fuerint oculati testes obseruationum Galilaearum, et illius doctrinam penitius imbiberint, illos siue hetruscos, siue Romanos, siue Bononienses, siue Genuenses et cæt. rogamus vt nobis satisfaciant, et quae \struck{hanc} hanc in rem parauerint ad eximium Geometram Genuensem D. Santininum mittant, quae fideliter ad nos transmitti curabit. Lutetiae Calendis Julii anni 1643
\note{« hanc » est écrit deux fois, la première biffée. La date est soulignée sur la page, avec un trait sous « 1643 ». « Santininum » : lu tel quel.}

\note{Sous la date, un diagramme occupe la moitié de la page : sept échelles de notes en escalier, chacune un faisceau de lignes horizontales entre deux montants, montant d'un degré de l'une à la suivante, et légendées au pied, de gauche à droite : « Hypodorien », « Hypophrigien », « Hypolidien », « Dorien », « Phrygien », « Lydien », « Mixolidien ». Chaque échelle porte les lettres de ses notes, de A (en bas de la première) à a et au-delà, avec le signe en forme de 4 déjà transcrit par le bécarre, et dans chacune trois lettres pointées placées à part, entre les deux montants (d, c, B ; D, C, B…), réunies par un trait vertical. Au niveau de la même ligne horizontale, chaque échelle porte en petit la syllabe de sa note : « mi la re », « re sol ut », « ut fa », « mi la », « la re sol », « sol fa ut », et à droite de la dernière un « mi » isolé ; c'est la « mese ou moyenne » dont parle la légende.}

visa vis de la mese ou moyenne du mode hypodorien l'on a le commencem\textsuperscript{t} des 6 autres modes ou harmonies, lesquelles estans commencées plus haut que l'A re dud\textsuperscript{t}. hypodorien suyuant les interualles marquez en cette table selon les Grecs, l'on entendra les 7 tons Et les 7 modes des anciens tous diuersifiez et chacun en leur lieu propre selon m\textsuperscript{r}. Doni, Et si l'on commence toutes ces six harmonies ou octaues au mesme ton de l'hypodorien, les modes seront tous au mesme ton hypodorien; les trois lettres interposées en chasque mode monstrent les \uncertain{propres} chordes de leur tetrachorde conioint./
\note{« mese » : lu « mize » ou « mese », le mot étant glosé par « moyenne ». « propres » : le début du mot est dans une ligne serrée, et n'est pas sûr.}

\end{document}
